% =============================================================================
% HIKET -- MANUSCRIPT DRAFT v1.  Written 2026-09-04 from HIKET_storyline_v3.tex.
%
% WHAT THIS IS. The first continuous-prose draft, on the architecture agreed in
% storyline v3: IMRaD outside, the narrative arc inside, Results and Discussion
% combined and nested at the subsection level (each subsection states a result in
% one paragraph and interprets it in the next).
%
% WHAT WAS CARRIED OVER from HIKET_main_manuscript.tex, and what was not:
%   CARRIED  Materials and methods, nearly wholesale -- it was current, checked
%            against the code on 2026-09-04, and is the most expensive material to
%            reproduce. Also the model-ladder framing (Yasso operational, the
%            simple models a robustness ladder), which was already correct there.
%   CARRIED  Two Discussion passages: the physical reading of the input
%            multiplier, and the candour passage about the level offset.
%   DROPPED  The whole Introduction (scaffolding notes only; rewritten from the
%            literature review completed 2026-09-04).
%   DROPPED  "The honest output is the spread" -- it claims the between-model
%            spread as a product, which was walked back on 2026-09-03.
%   DROPPED  The locally stratified calibration strand (out of this publication).
%
% STYLE RULES for this draft, set by Lorenzo 2026-09-04, and they are binding:
%   * No jargon. Where a technical term cannot be avoided, define it in plain
%     words at first use, then use it consistently.
%   * No implicit pronouns. Do not open a sentence with "This", "It" or "These"
%     standing for something in the previous sentence. Name the thing again.
%   * Be redundant rather than compact. Repeat the subject. A few more words is
%     the right trade.
%   * Scientific prose should be boring. Old information first in the sentence,
%     the new information at the end.
%
% NARRATIVE FRAMEWORKS IN USE (see memory narrative-tools-ocar-jung).
%   OCAR (Schimel) -- the public arc, now INTERNAL to the IMRaD shell: Opening and
%     Challenge live in the Introduction, Action in Materials and methods, Resolution
%     across Results and discussion. Signposted in source comments only.
%   Gopen & Swan -- sentence level, applied throughout: old information first, the
%     payload in the STRESS position at the end of the sentence.
%   ABT -- diagnostic only, used sparingly. Kishotenketsu -- still a candidate.
%
%   Jungian archetypes -- HIDDEN. A private map only; these words must NEVER appear in
%   the text, or the effect is destroyed. RESCUED HERE 2026-09-07 from the headers of
%   HIKET_main_manuscript.tex and HIKET_storyline_v2.tex, both of which were deleted as
%   superseded on that date; this draft is now the map's only working home.
%     Protagonist  = the INITIALISATION PROBLEM -- how to start a non-equilibrium
%                    simulation with no history -- with transient init as the move.
%                    Held constant throughout.
%     Shadow       = the observed accumulation that bends over, which is precisely what
%                    a steady-state start CANNOT produce, because an equilibrium start
%                    is already saturated. Lives in the Introduction's Challenge.
%     Trickster    = sigma_input, the shape-shifter that looks like a nuisance
%                    parameter and turns out to carry the finding.
%     Mentor       = the published parameterisations and the group's own lineage: real
%                    but boreal-blind, then freed. Treated collegially, NEVER as an
%                    antagonist -- Aleksi and Mikko are coauthors.
%     Persona      = good R-hat and physical-looking bias, masking artefacts.
%     Unconscious  = the unrecoverable pre-observation history; the humus pool that
%                    remembers equilibrium, deep time made physical.
%     Self         = the fair, equal-information, exactly-integrated, homogenised frame.
%     Boon         = a method and a relocated question, not a verdict on Finland. The
%                    forward experiment is the return with it: the displacement made
%                    consequential, and bounded.
%
% NUMBERS. All results are from the production run 20260831_1624* (arm B).
% =============================================================================
\documentclass[11pt]{article}
\usepackage[utf8]{inputenc}\usepackage[T1]{fontenc}
\usepackage[margin=2.5cm]{geometry}
\usepackage{graphicx}\usepackage{amsmath}\usepackage{amssymb}
\usepackage[hidelinks]{hyperref}
\usepackage[round]{natbib}
% Parent for the standalone appendices in appendices/*.tex.
\usepackage{subfiles}
\usepackage{caption}\captionsetup{font=small,labelfont=bf}
\usepackage{xcolor}
\usepackage{titlesec}
\titlespacing*{\paragraph}{0pt}{0.8\baselineskip}{0.4em}
\linespread{1.08}

% marks a place where the draft is not finished; nothing else should be coloured
\newcommand{\gap}[1]{\textcolor{red!70!black}{\small[\,#1\,]}}
\graphicspath{{figures/}}
% One figure macro, so width and placement stay consistent across the paper.
\newcommand{\hfig}[4]{%   #1 file  #2 width  #3 label  #4 caption
  \begin{figure}[htbp]\centering
    \includegraphics[width=#2\linewidth]{#1}
    \caption{#4}\label{fig:#3}
  \end{figure}}
% A figure held for the supplement; the name is kept so the file can be found.
\newcommand{\suppfig}[1]{\textcolor{blue!55!black}{\small[Supplementary figure: \textsf{#1}]}}
\newcommand{\figref}[1]{\textcolor{blue!55!black}{\small[Figure: \textsf{#1}]}}

% ---------------------------------------------------------------------------
% TITLE. Chosen 2026-09-11 with Lorenzo, after two rounds. The alternatives are
% kept below so that swapping is one edit.
%
% WHY THIS ONE. The previous title ("what four decades of Finnish soil measurements
% require of six forest carbon models") had two faults. It billed the six models,
% when three of them are benchmarks and do not deserve that share. And it PROMISED a
% finding ("what ... require of") without stating one, which is the fault a title in
% this literature can least afford.
%
% The title now states the finding and the mechanism in one declarative clause, and
% stops. Endings that promise a consequence -- "with consequences for projection" --
% were tried and rejected: a clause any paper could append. The consequences live in
% the abstract.
%
% ⚠ THE VERB IS "DISPLACES ALONG", NOT "REVEALS". Relaxing the steady-state
% assumption does NOT reveal the trade-off: the trade-off exists under either
% initialisation, and a transient start only NARROWS it, partially (transit time and
% the input multiplier still correlate -0.64 to -0.73). That correction was made on
% 2026-09-03 and must not be reverted through the title. "One displacement along a
% trade-off" is the Conclusions' own wording.
%
% ⚠ "TURNOVER", not "transit time" -- deliberate. Sierra et al. 2017 explicitly adopt
% turnover time as stock over flux, so the plain word is not the discouraged term, and
% it reads to a wider audience. The Methods do the precise naming.
%
% ⚠ TURNOVER IS NOT IN THE TITLE AS A DIRECTION, and must not be. The input
% displacement is universal (effective flux 3.19-4.34 against 2.511, all six models);
% the turnover displacement is GRADED -- ratios 0.73 / 0.89 / 1.03, and Yasso20 is not
% displaced at all. Yasso20 is the negative control. A title claiming "faster turnover"
% would contradict it.
%
% ⚠ FINLAND IS NOT IN THE TITLE, by the rule that Finland is the case and not the
% claim. The abstract and the first lines of the Introduction say where the
% measurements come from.
%
% ALTERNATIVES, in the order they were considered:
%  (b) ... along a trade-off that the measurements narrow but do not resolve
%      [same stem, with the contribution AND its limit in the tail]
%  (c) ... displaces them towards a larger litter input, along a trade-off the
%      measurements narrow but do not resolve   [adds the direction; 23 words]
%  (d) What a transient initial state can and cannot settle: litter input, turnover,
%      and projection in forest soil carbon models
%  (e) Soils that are still catching up: relaxing the steady-state assumption in
%      forest soil carbon models
%  (f) Forest soils are not where models start them: consequences of estimating
%      the initial state of soil carbon in Finland
% ---------------------------------------------------------------------------
\title{\vspace{-1.2cm}Relaxing the steady-state assumption displaces\\
forest soil carbon models along the trade-off\\
between litter input and turnover}
% Author list as of 2026-09-15 (Lorenzo). Order not yet decided; affiliations pending.
% Boris Tupek (litter product) and Olli-Pekka Tikkasalo added 2026-09-15.
% Mikko Peltoniemi REMOVED for now -- Lorenzo has not yet asked him; may return.
\author{Lorenzo Menichetti, Aleksi Lehtonen, Samuli Launiainen, Boris {\v T}upek,
        Olli-Pekka Tikkasalo, \gap{possibly others; order and affiliations to be decided;
        Mikko Peltoniemi to be asked}}
\date{Draft of \today}

\begin{document}
\maketitle

\begin{abstract}
% ROUND 5, 2026-09-23: first draft of the abstract (~290 words). Numbers: effective
% flux 3.19-4.34 / J_bar 2.511 = 1.27-1.73; MTT ratios 0.73 / 0.89 / 1.03; headroom and
% sinks from manuscript/figures/forward_scenarios.rds.
Soil carbon models used in national greenhouse gas inventories are commonly started
from a soil in balance with its litter input, although the soils of managed forests are
rarely in balance. We calibrated six soil carbon models, the three versions of Yasso,
the model used in the Finnish inventory, and three simpler benchmark models with one to
three pools, against soil carbon measured in three campaigns between 1986 and 2024 on
$456$ permanent plots of the Finnish National Forest Inventory. Each model started from a
soil in balance in 1917 and was run forward along a reconstructed history of litter
input. The litter input and the input of 1917 were estimated together with the model
parameters in a Bayesian calibration, with an error model in which observations share
part of their error by plot, by region and by campaign. All six models represented the
national trajectory of soil carbon, and none could be distinguished from the others by
its fit. The calibration placed the models at a litter input $1.3$ to $1.7$ times the
input supplied by a tree-litter product and, in Yasso07 and Yasso15, at mean transit
times $27$ and $11$ per cent shorter than published, whereas Yasso20 was not displaced.
Litter input and transit time trade off against each other, and returning to the
published transit times cost little in fit, so the observations suggest a direction
rather than impose a value. The position on this trade-off changed the projection of
Yasso07 only. Under a constant litter input the models placed the soil $26$ to $42$ per
cent below its state of balance, consistent with an earlier reconstruction, while the
projected sink fell by a third to a half within $60$ years; the size of that sink depended
on how depleted the soil was estimated to be at the start. None of the models
represented the differences between individual plots, and about half of what the models
left unexplained could be recovered from site variables that already exist. The study
provides a baseline for further development of the models used in forest soil carbon
accounting.
\gap{Check the length against the journal's limit, and revisit once the Limitations and
the discussion of the starting state are settled.}
\end{abstract}

% =============================================================================
% =============================================================================
% INTRODUCTION -- RESTRUCTURED 2026-09-07, in answer to Lorenzo's 20 notes on the
% first draft. The old prose is preserved verbatim in the comment block at the end
% of this section, so nothing was removed without a record.
%
% WHY IT WAS RESTRUCTURED. The subject of the paper -- the initial state -- first
% appeared halfway down page 2, inside a paragraph about uncertainty percentages,
% introduced by "The reporting chain also...". The Challenge therefore arrived as
% the fourth item in an unlabelled list rather than as one named thing. Two faults
% followed: the argument was made twice (the theory in the old sec.4 explained an
% assumption the reader had met a page earlier), and meta-commentary about what the
% paper is NOT was standing in for the argument.
%
% NEW ORDER (2026-09-07). 1 Opening, biology first, then the campaigns, then why the
% rate and not the stock. 2 The soil term is modelled, and where its uncertainty comes
% from. 3 The convention. 4 The two consequences. 5 Why Finland is the wrong place for
% the assumption. 6 Objectives.
%
% RE-THREADED 2026-09-15 (Lorenzo: "everything is a bit all over the place" -- the
% Finnish inventory was introduced TWICE, as a per-hectare measured rate in sec.1 and
% as a Mt CO2 balance in sec.2). The thread now runs general -> particular in one
% direction only:
%   1 The general problem (unchanged, without the campaign numbers).
%   2 How the soil term is obtained: other countries, then Finland's chain.
%   3 NEW: what the Finnish chain has reported, and WHY Yasso reports a sink (the
%     NID's own mechanism: spin-up on NFI6 litter, litter risen since).
%   4 The measurements, and the new chance the third campaign gives (moved here
%     from sec.1, with Figure 1).
%   5 The convention, opened by the "started somewhere" paragraph, then as before.
%
% DECIDED WITH LORENZO, 2026-09-07:
%   * Opening move = biology-anchored ("a soil that is still catching up with its
%     forest"), NOT the policy framing. Per the weight budget: anchored in biology.
%   * The review of national soil-budgeting methods stays at ONE sentence here and
%     moves to the Discussion, where it attaches to what the result means for other
%     inventories. See the \gap marker at the end of Results and discussion.
%
% DISPOSITION OF THE 20 NOTES. [R] resolved here, [O] still open, your call.
% [R]  1 Opening too method-specific -> new generic opening; campaigns come third.
% [R]  2 Qualify the national inventory effort -> NFI / BioSoil / 2024 each named.
% [R]  3 "connect these two with a but" -> the two rates are now one sentence.
% [R]  4 CI-includes-zero "seems to mean nothing" -> KEPT, and now doing visible
%        work: the levelling-off IS the contrast between the two rates, so the
%        paragraph states which rate is robust (9% offset) and which is fragile
%        (2%). Hiding it would leave the hook resting on the weaker number.
% [R]  5 Reporting basis is too much M&M -> demoted to a footnote.
% [R]  6 State vs trajectory, and the budget is a rate not a state -> new closing
%        paragraph of sec.1, phrased as you put it: the inventory reports the
%        annual change, so the trajectory is the quantity it needs.
% [O]  7 "as many other countries", short review of national budgeting methods ->
%        ONE sentence here; the review moves to the Discussion (decided). The gap
%        marker is in place; the primary sources for three of the rows are not yet
%        read, so the review cannot be written until they are.
% [R]  8 "uncertain enough to matter" was cool but doubtful -> sentence deleted.
% [R]  9 "when the national balance is close to zero, meaning what?" -> replaced;
%        the sector balance is now named explicitly with its figure.
% [R] 10 The sign argument is self-evident/circular -> replaced. The point is no
%        longer size x uncertainty but WHERE the uncertainty comes from: the
%        reported quantity is a rate, the rate is modelled, and one input to the
%        model is never measured anywhere in the chain.
% [R] 11 "The reporting chain also..." -- also with respect to what? -> the whole
%        move is gone. The initial state is now introduced as the subject, in its
%        own subsection, not as an afterthought in a list.
% [R] 12 "steady state" is standard terminology, and "run into a steady state"
%        reads as mockery -> "steady-state initialisation" used throughout, and
%        the Lehtonen 2016 quote is now given as a premise, not as a specimen.
% [R] 13 ROMULv appears from nowhere -> introduced in the same sentence.
% [O] 14 CBM-CFS3: "is this so radically different?" -> largely CONCEDED. A cycle
%        oscillating about a mean is still a balance, so "instructive exception"
%        was too strong; it is now "a partial exception", and the difference that
%        survives is that the last disturbance is dated. Check whether even that
%        is worth the clause.
% [R] 15 Liski is the one conceptually close antecedent -> moved up beside the
%        other criticism, where it now carries the novelty claim, and removed from
%        its second appearance so the point is made once.
% [R] 16 "the criticism is not what this paper contributes" is too open about our
%        intentions; the elephant problem -> all three meta-sentences deleted.
% [R] 17 Wutzler makes little sense as written; does it involve fitting? -> YES,
%        and you were right. Rewritten: the practice is choosing parameters so the
%        modelled steady state matches the observed stock; a still-accumulating
%        soil then implies a very large equilibrium stock, so the fit overestimates
%        the decay rate of the slowest pool. Source: literature/02, verified.
% [O] 18 "Wherever the assumption has been relaxed... Really? And Liski before?" ->
%        the sweeping quantifier is gone. The claim is now the narrow one: nobody
%        has replaced the assumed state with one ESTIMATED FROM A MEASURED CHANGE
%        in the same soil. Liski ran the transient simulation and assumed the 1922
%        state, so Liski is antecedent for the run, not for the state. If you want
%        "dominant" anywhere, we still owe a count.
% [R] 19 The theory is too disperse and in the wrong place -> the definition moved
%        up into sec.3 as you asked, and the consequences now have sec.4 to
%        themselves, at half the length.
% [R] 20 "turns on both" unclear; what does the product mean; is it a trade-off ->
%        plain language throughout: "the results depend on both"; the product is
%        spelled out as what a stock can and cannot identify; "trade-off" replaced
%        by "exchange", which is what it is.
%
% SECOND PASS, same day: both subsection titles claimed less than the argument does.
% [R] 21 "The soil term of the national inventory is produced by a model" read as a
%        fact about Finland -> now "The soil term of A national inventory is modelled
%        rather than measured", and the prose runs generic first (any country
%        reporting in detail reports a modelled quantity) with Finland as the
%        instance. The subsection now ends on a hook into sec.3: a model asked to do
%        that has to be started somewhere. Sec.3 answers it in the same terms -- the
%        model families used for national reporting are the ones started at steady
%        state -- so the two subsections are joined by the argument rather than by
%        adjacency, which is the link you asked for.
% [R] 22 "Finnish forests are a managed landscape and not an equilibrium one" implied
%        the problem is Finnish -> now "A managed forest landscape is not an
%        equilibrium landscape", opening on the general claim (a forest whose biomass
%        changed within the residence time of soil carbon leaves a soil out of
%        balance, and such landscapes are common) with Finland as a well-documented
%        case. The generalisation is the honest one: what is special about Finland is
%        not the disequilibrium but the fact that it has been measured three times.
%
% =============================================================================
% MATERIALS AND METHODS -- ROUND 1 of your notes, 2026-09-11. Handled so far:
%
% THE ERROR MODEL. Three defects, all found while checking your notes, none of
% which moves a number (production is arm B: HIKET_CORRELATED_LIK=1 makes
% hiket_total_sigma() return HIKET_SIGMA_TOT=0.800 before sigma_obs_fixed is ever
% reached, and the predictive stage reads sigma_total from the run's metadata).
%   1. "the measurement error of the soil carbon determination alone, which is
%      0.43" was NOT a measurement error. The code computes
%      sigma_obs_fixed <- sd(SOC_obs)/mean(SOC_obs), i.e. the BETWEEN-PLOT CV of
%      the observed stocks (0.44/0.47/0.53 by campaign) -- the spread of soil
%      carbon across Finland, not a determination error. The comparison is DELETED;
%      we have no sourced determination error (checked Lehtonen & Heikkinen 2015 --
%      their budget is biomass and litter, no SOC measurement enters that chain).
%   2. 0.800 was described as "the measured spread of the model residuals". The
%      measured spread is 0.71-0.74 and 0.72 was the plug-in value; 0.800 is a
%      DELIBERATE step above it. Now said that way.
%   3. The four numbers were called variances. They are standard deviations:
%      sqrt(0.117^2+0.396^2+0.030^2+0.685^2) = 0.8004, the fixed total. As
%      variances they would imply a total SD of 1.12. Fixed, and the quadrature
%      is now stated. Also repaired a sentence that had lost its opening clause.
%
% MEAN TRANSIT TIME (your note at "The six models"). Sierra et al. 2017 read in
% full; PDF in literature/09_transit_time/. The change is a RENAME, not a
% recomputation -- their Eqn (2) is our formula character for character:
%     MTT = (1,...,1) (-A)^-1 I/sum(I)
% and their Eqn (4), MTT = sum(x_ss)/sum(I), is our numerical route (unit input,
% run to balance, sum the pools). Eqn (4) also ANSWERS YOUR UNITS COMMENT: the
% quantity is a stock divided by a flux, so it is in years; ours reads as a bare
% stock only because sum(I)=1 makes the two numbers coincide. Now written as the
% ratio. Verbatim from the paper: "All these different uses of the term residence
% time make it difficult to unambiguously apply it in carbon cycle research. We
% therefore do not consider this term any further in this manuscript and
% discourage its use in further research"; and "In autonomous multipool systems at
% steady state, the turnover time is equivalent to the mean transit time" -- which
% is exactly our construction (model frozen at the reference climate, at steady
% state), so a clause now says so, and the Methods state that we report the
% transit time of the FROZEN system and not of the real soil.
%   * Named once as "mean residence time" in the Methods, deliberately, because
%     that is what most readers know it as. Do not delete that sentence.
%   * "turnover time" gets one clause for the same reason.
%   * MRT -> MTT in figure labels and in T_A1's column header (MTT is Sierra's own
%     abbreviation). 27 occurrences renamed in the prose.
%   * The fine-root sentence in "What might be missing" said "mean residence time"
%     for ROOT LIFESPAN -- a different quantity, from a different literature. It now
%     says "mean lifespan" rather than being renamed to transit time.
%   * CODE IS NOT RENAMED (decided): intrinsic_mrt.R, the mrt columns and the
%     build_*_mrt_*.R filenames stay, because the figure staleness guards key on
%     them. Recorded in RELEASE_NOTES.md, which is the new home for
%     packaging-for-publication items; CLAUDE.md points at it.
%
% SECOND PASS, same day -- ALL of the remaining notes are now written:
%   * workbook NOT named and NOT cited -- see the note in the subsection itself.
%   * STONINESS, your "I fear we don't": WE DO. The workbook's stocks are already
%     corrected for bulk density and coarse fragments. The paragraph now says what we
%     take rather than what an earlier assembly omitted.
%   * "EXACTLY": re-ran build_soc_homogenized.R. On the official calculation's own
%     exclusions (n=446) it returns 59055 and 61047 kg/ha against targets of 59055 and
%     61047. The text states the check, the basis and the n, and quotes agreement to
%     0.01 t/ha rather than asserting "exactly".
%   * "which earlier assembly" -> that comparison is GONE.
%   * layers: written as standardisation, in your words. Sampling years: written as a
%     choice, not a correction.
%   * litter first year: short and neutral; the biomass-differencing MECHANISM is no
%     longer asserted, only the value and the treatment.
%   * 1917 = Finnish independence, said in one clause, with the point that nothing
%     depends on the exact year.
%   * PRESCRIBED SHAPE: now three paragraphs, because you were right that it is one of
%     the load-bearing assumptions. It is NOT the growing-stock ramp -- the default
%     changed on 2026-08-20 to the Liski et al. 2006 reconstructed input to soil. The
%     text gives the three decisions taken in adopting it (tree-derived terms only, to
%     match the post-1985 product; per-hectare using their own implied upland area;
%     shape only, not level), why standing volume is the wrong driver (a stock, not a
%     flux, and it misses harvest residue, so it holds the input at its 1917 minimum
%     through the heaviest-cutting decades while the reconstruction peaks ~1960), and
%     the one residual bias, stated not corrected, which runs in the conservative
%     direction. Source: Model_functions_real_data_transient/preinit_input_shape.R.
%   * WHY YASSO'S RATES ARE FIXED: your reason, written out -- rates and fractions are
%     strongly equifinal so fitting both buys little, and of the two the rate is the
%     more fundamental because it defines the pool's position on the decomposability
%     spectrum by setting its mean transit time.
%   * WHY THE SIMPLE MODELS ARE PINNED, NOT FIXED, TO ICBM: your reason -- ICBM's rates
%     come from an arable experiment, so fixing there would impose a displaced value,
%     while fitting freely would give the simple models freedom Yasso does not have.
%     The asymmetry is admitted and the posterior width is reported as the diagnostic.
%   * INTEGRATION: "exactly rather than approximately" OVERSTATED THE YASSO CASE.
%     matrixexp() in the Fortran is scaling-and-squaring with a 20-term Taylor series --
%     numerical, though of the matrix exponential and not a time-stepping solver. The
%     text now draws the distinction that actually matters: exact solution of the linear
%     system over the step, versus an explicit step of the derivative.
%   * SAMPLER: ter Braak & Vrugt 2008 (DEzs) via BayesianTools 0.1.8 under R 4.3.1, and
%     the three internal chains per run are stated, so five runs give fifteen chains.
%   * Two \gap markers also closed: the sampling-year groups (geographic, with the
%     decisive 2006/2024 control) and the 1985 down-weighting (replaced by the campaign
%     offset, not stacked with it).
% STILL OPEN in M&M: the weather product's provenance (asked), and the parameter-class
% by prior-criteria table.
%
% =============================================================================
% ROUND 2 -- 16 marginal notes on the restructured Introduction, 2026-09-11.
% All 16 acted on. Order of the section changed once more, as you asked in note 16:
%   1 the soil catching up -> 2 the term is modelled -> 3 the convention -> 4 why a
%   managed landscape is not in balance -> 5 what a balanced start cannot do -> 6 aims.
% Each subsection now answers the question the previous one raised.
%
% [R]  1 "holds the carbon of every tree" is absolute and wrong -> "Part of the carbon
%        held in a forest soil came from trees that died long before the trees now
%        standing on the site." Same move, no totality claim.
% [R]  2 "store" reads like a shop -> "the slowest fraction".
% [R]  3 "slowest ... slowly" -> repetition removed.
% [R]  4 Le Noe: most SOC models ARE built on observations -> claim narrowed to what
%        the review actually reports, namely testing against a time series that was
%        NOT also used to build the model. Note the trap we avoid by demoting it: if
%        we lean on "nobody validates on time series", the next question is whether WE
%        provide that, and we calibrate with a spatial holdout, not a temporal one.
% [R]  5 "one of the few countries" is contestable (Sweden) -> comparative claim
%        deleted. Replaced by the functional one: estimating a rate, or a state from
%        it, requires the same soil sampled more than once, and Finland's plots have
%        been sampled three times. Nothing there to argue with.
% [R]  6 Abbreviations -> C and SOC introduced at first use and used throughout. Kept
%        "soil carbon" in the two places where the quantity includes dead organic
%        matter (the inventory reports DOM and SOM aggregated), so that we do not
%        mislabel the inventory's own aggregate as SOC.
% [R]  7 "the whole soil profile", specifically? -> footnote now says it: organic
%        layer, mineral soil to 40 cm, modelled extrapolation below, integrated to one
%        metre OR to the depth the auger reached, whichever is shallower. Answers the
%        bedrock objection in the same breath.
% [R]  8 The falling sink is the pitch -> the 84 per cent decline and the sector
%        turning into a source now OPEN the paragraph and carry it.
% [R]  9 Uncertainty is not the load-bearing point -> demoted from its own paragraph
%        to a clause on the soil term, where it belongs.
% [R] 10 Steady state is often analytical, and the assumption is the point, not the
%        implementation -> definition rewritten: solving the equations or running the
%        model forward is "a matter of implementation"; the assumption is named as
%        what matters.
% [R] 11 ROMULv is not in the inventory -> CORRECT, and the draft was wrong. Verified:
%        the NID says "For Forest Land Remaining Forest Land, the Yasso07 model was
%        applied" (FI_NID_2024.txt:15977), Yasso07 alone. Lehtonen et al. 2016 is a
%        nationwide EVALUATION of ROMULv and Yasso07, not the reporting instrument.
%        The quote is kept, the attribution fixed: "The nationwide assessment that
%        evaluated Yasso07 and ROMULv against Finnish soil measurements".
% [R] 12 CBM-CFS3 is an apparent exception and really a confirmation -> rewritten as
%        exactly that: a state cycled to a repeating average is still a state in
%        balance, and what the procedure adds is a dated last disturbance rather than
%        a history of directional change. Caveat: Kurz et al. 2009 is still read only
%        at second hand, so the claim is kept to what the abstract supports.
% [R] 13 Liski deserves a structured statement -> own paragraph, built as you asked:
%        first everything the two studies share (country, period, model family, NFI
%        source, and the inventory still follows their method), then the single
%        assumption that separates them, then what we do instead. Framed as lineage.
% [R] 14 Equifinality is generic, not a consequence of steady state -> sec.5 now opens
%        with the general trade-off between input and residence time, in three
%        sentences, and steady state follows as its extreme case in one: in balance
%        the stock IS the product, C* = J/k, so the product is all a stock can fix.
%        This also stops the Introduction from contradicting the Results, where the
%        ridge is measured under a TRANSIENT start and survives (-0.64 to -0.73).
% [R] 15 The trajectory consequence is the main one -> given the second paragraph of
%        sec.5 and named as the one the paper turns on.
% [R] 16 Move the managed-landscape argument up, to answer "why depart from steady
%        state" -> done; it is now sec.4, directly after the convention and Liski,
%        and the consequences follow it as sec.5, adjacent to the objectives.
%
% NEW MATERIAL, and the reason this round took a download: Peltoniemi et al. 2006
% (FEM 232:75-85) is the single citation the inventory rests its initialisation on,
% and it had never been read. It is now in literature/02_steady_state_initialisation/.
% What it says, and why the new paragraph in sec.3 is fair rather than adversarial:
%   * the numbers are a PRECISION decomposition -- the initial state's share of the
%     SD of the annual sink: 77% (1989), 30% (1998), 21% (2004). "Cancel out" in the
%     NID's paraphrase means "no longer dominant", not "negligible".
%   * the initial state they perturbed WAS ITSELF A STEADY STATE ("At the start of the
%     calculation period (1988), the soil model was in steady state with the input and
%     climate of the first year"; the Appendix derives it "based on inputs and model
%     parameters"). So they propagated the uncertainty OF an equilibrium value, never
%     the uncertainty OF the assumption of equilibrium.
%   * they say so themselves: "our results represent the precision, not accuracy".
%   So we never contradict Mikko. We say the inventory uses a precision result to set
%   aside a question about accuracy. He is a coauthor; this framing respects that.
% HELD BACK for the Discussion, deliberately, both from the same paper:
%   * their conclusion is conditional -- initialisation matters most "if there are no
%     long-term records of data on forests prior to the years for which the soil sinks
%     are estimated". We DO build such a record (1917-1985) and the state still moves,
%     which is a result, not a premise.
%   * litter uncertainty was large but autocorrelated at 0.99, so it did not propagate
%     into annual sinks; at 0.9 the soil sink SD went 2.6 -> 3.9 Tg, at 0.7 -> 5.5 Tg.
%     That is our correlated-likelihood argument, from the same group, in 2006.
% =============================================================================
\section{Introduction}

\subsection{A soil that is still catching up with its forest}

Forest soils are the largest carbon (C) store of the forest ecosystem: in the boreal
zone, more of the C held by a forest lies in its soil than in its trees
\citep{Pan2011}. 

But what is most crucial, more than the snapshot of total C stocks, is the C balance: if soils are gaining or losing C. A national greenhouse gas inventory reports the annual change in the stock rather than
the stock itself. 

A soil that is still gaining C and a soil that has stopped gaining C (wich is defined as equilibrium, since there is still fluxes in and fluxes out that just balance each other) can hold the same
amount of C, and the two soils differ only in where they have come from and where they
are going. The trajectory is therefore the quantity that decides whether a
reported sink will persist, and for how long.

Such trajectory is influenced by the history of the soils considered, even land use changes relatively far in time.
While litter falls to the soil continuously and most of its C
returns to the atmosphere within a few years, a small but consistent fraction of that C joins
material that decomposes on a scale of decades to centuries.
Because this slower C persists for so long, a forest whose management changed even a century ago
can stand on a soil that is still adjusting to such change.
Such a soil, not in equilibrium, is gaining or losing C at a certain rate.


\subsection{The soil term of a national inventory is often modelled rather than measured}
Countries determine the soil organic C (SOC) in their inventories in more than one way. Some
measure it. Sweden estimates the change in the C stock of its forest mineral soils
from a soil inventory carried on the permanent plots of its national forest inventory
since 1993, on a ten-year cycle, and reports the measured difference
\citep{SwedenNID2026}. Germany does the same from two national forest soil inventories
two decades apart \citep{Gruneberg2014}. Repeating a national soil survey often
enough to yield a stock change is expensive, however, and the more common route is a
decomposition model driven by the forest inventory: Finland, Norway, Switzerland and
Austria run Yasso07 on their inventory plots \citep{StatFin2025NID,NorwayNID2025,Didion2016},
and Canada runs CBM-CFS3 \citep{Kurz2009}. The modelling route inherits a difficulty
that Norway's inventory states plainly, that ``in a managed forest landscape neither
forest biomass, soil nor DOM pools can be assumed to be in an equilibrium state''
\citep{NorwayNID2025}. And yet any model needs to start from a certain state, and usually that state is determined by assuming equilibrium conditions at some point past in time. This means assuming a certain mean input level as constant, and calculating the C stocks from that assumption, before starting a transient simulation.

Finland, specifically, reports to the United Nations Framework Convention on Climate Change through a chain in
which forest inventory data and biomass models supply an estimate of the litter falling
to the soil, and the decomposition model Yasso07 converts that litter estimate into an
annual change in the soil C stock, for dead organic matter and soil organic matter
combined to a depth of one metre \citep{StatFin2024NID}, following the method of
\citet{Liski2006}. No measurement of soil C enters the Finnish chain at any point.

\subsection{What the Finnish chain has reported, and why}
The quantity produced in that way has been, for most of the reporting period, a sink, and
a large one. The combined pool of dead organic matter and soil organic matter on
mineral forest land was reported as a sink of about $10$ million tonnes of carbon
dioxide equivalent per year at the beginning of the 2000s \citep{Luke2025news}, and
in the 2024 submission it was still a sink of $4.8$ million tonnes in 2022, against a
living biomass sink of $12.6$ million tonnes \citep{StatFin2024NID}. The soil term
carries a reported uncertainty of $31.5$ per cent against $16.3$ per cent for the
biomass term \citep{LehtonenHeikkinen2015}. The resulting net balance of the whole land-use (LULUCF) sector is now
close to zero: the C sink of Finnish forest land has fallen by more than $80$ per cent
since 1990, and the land use sector as a whole has become a net source
\citep{StatFin2025NID}. A term of the size and uncertainty of SOC can decide the sign of
such a balance.

In the 2025
submission, which recalculated the litter input from the thirteenth National Forest
Inventory, the mineral soils of forest land turned ``from a net sink to a net source
from 2021 onwards'', emitting $0.4$ million tonnes of carbon dioxide in 2023
\citep{StatFin2025NID}. The reason given is the litter input: the model is started from
a soil in balance with the litter input of the sixth National Forest Inventory
(1971--1976); such input rose for decades with the growth of the standing biomass
and with the residues of an increasing volume of fellings, and the modelled soil
accumulated C in response; the recalculated litter input then fell, mainly through a
lower estimate of foliage biomass, and the modelled soil began to lose C
\citep{StatFin2025NID}. The reported soil term is therefore, in both directions, the
response of a modelled soil to a litter input that departs from the level the soil was
started in balance with. Whether the real soil was in balance with its litter input in
the 1970s, and how far it has come since, is not something the chain can tell, also because there are
no soil measurement to inform the model and everything follows from the mode3ling assumptions.
% Per-hectare reading (literature/10_national_soil_terms/README.md sec.3): mineral forest
% land 15.9-16.0 Mha (NID 2025 Table 6.4-1) => ~10 Mt CO2 eq ~ +0.17 tC/ha/yr (early
% 2000s), 4.8 Mt ~ +0.08 (2022, 2024 submission), 0.4 Mt source ~ -0.007 (2023). Against
% the measured +0.399 (1989-2006), +0.117 (2006-2024), +0.259 (full). Same direction,
% larger swing in the model. Bases differ (area-weighted national, incl. dead wood, vs
% the balanced 310-plot soil set) -- order of magnitude only. Not quoted in the text yet.
\gap{one clause or one panel: the reported term per hectare against the measured
rate, on the CRT series --- decide after the colleagues confirm the $-10$ figure}

\subsection{Three campaigns on the same plots}
But the permanent plots of the Finnish National Forest
Inventory have been sampled now three times: a first campaign whose fieldwork ran from 1986
to 1995, the European BioSoil programme in 2006, and a third campaign in 2024.  On the
plots measured in all three campaigns, the mean soil organic carbon (SOC) stock rose
from $64.5$ tonnes of C per hectare to $71.5$ and then to $73.6$, a gain of $7.0$ tonnes
over the first interval and a gain of $2.1$ tonnes over the second.
Three points give already a trajectory, that can be used to constrain a SOC model.

\hfig{F3_mean_SOCalt}{0.86}{observed}{\textbf{The measured soil carbon stock of Finnish forest plots in three campaigns.} Each point is the mean stock across the plots measured in all three campaigns ($n=310$), and the bars show the $95\%$ confidence interval of that mean. The first campaign is drawn at the mean year in which its plots were actually sampled, and the bar low in the panel shows that the sampling of the first campaign ran from 1986 to 1995. Each segment is labelled with the rate of change over that interval and with the $95\%$ confidence interval of the rate, calculated from the paired difference between campaigns on the same plots. The stock rose over both intervals, and the rise over the second interval was smaller than the rise over the first.}

The measured soil has been gaining C, as the inventory reports, and the gain has been recently
slowing down.



\subsection{Steady state as a starting condition}

Any SOC model needs to be initialized somewhere. The model has to be told how much C the soil already holds on the day the
simulation begins, and how that C is divided between fractions which decompose at
different speeds, for a soil whose history was never observed. 

The standard solution is steady-state initialisation: the initial state is taken to be
the state at which the C stock would stabilize (zero change) under a litter input and a climate
held fixed. That state can be obtained by solving the model equations directly or by
running the model forward until the stock settles. The choice between the two is a
matter of implementation and what matters is the core assumption: that the
SOC was in balance with its litter input on the day the simulation began.

Most model families used for national reporting of forest soil C are initialized in that way. The
nationwide assessment that evaluated Yasso07 and ROMULv against Finnish soil measurements
initialised both models at steady state \citep{Lehtonen2016}. The spin-up
procedures of global land models rest on the same assumption \citep{Luo2016}, as do
applications of the model RothC \citep{ColemanJenkinson1996,Contreras2026}. The Canadian
forest C model CBM-CFS3 looks at first like an exception, because the model builds its
starting state from repeated cycles of growth and disturbance \citep{Kurz2009}; but a cycle around a constant average is still representing a trajectory that can be approximated as flat, and cannot represent any other trajectory.

In the Finnish inventory, the initial state is obtained by running the model under an unchanging litter
input until the C stock stops changing, using the average litter input of the sixth
National Forest Inventory (1971--1976) and the average weather of 1960--1990, and the
justification offered is that ``approximately ten years of simulation since spin-up is
enough to cancel out the effect of the spin-up level''
\citep[\S6.4, citing][]{StatFin2024NID,Peltoniemi2006}.

There is already considerable uncertainty around that assumption. \citet{Peltoniemi2006} repeated the Finnish calculation as a Monte Carlo
experiment and found the uncertainty of the soil C sink to be ``dominated by the soil
model initialization, but the effect decreased with time'': the initial state contributed
$77$ per cent of the standard deviation of the annual sink in the first year, $30$ per
cent after ten years, and $21$ per cent after sixteen. 
But the initial state whose uncertainty was propagated was itself a steady state, so what the experiment measured is the spread
around an assumed level rather than the full uncertainty in the trajectory.

\citet{Wutzler2007}, working
with the Yasso model, examined the common practice of choosing model parameters so that
the modelled steady state matches an observed C stock. A soil that is in fact still
accumulating implies an equilibrium stock far above the stock observed, so fitting under
the assumption of balance systematically overestimates the decay rate of the slowest C
pool, and overestimates the stock of a recently disturbed site. The same conclusion was
drawn, in the same year, in a review of seven soil C models for national accounting written
largely by the group that built the Finnish chain \citep{Peltoniemi2007}: the assumption of a
soil in balance with its current inputs ``is likely to be violated in most applications'', and
a model whose parameters are set so that its equilibrium matches a measured stock needs ``a
further correction of the parameters and the pools'' if the accumulation over a century is not
to be underestimated.
The steady state assumption biases the parameters retrieved by model--data fusion and not only
the initial state \citep{Carvalhais2008}, and later work has measured what the choice of
initial state costs in predictive error \citep{Lee2020,Kanari2022}.

The closest predecessor of the present study is \citet{Liski2006}, and the two studies
share the conceptual approach of pushing the steady state assumption as far back in time as possible. Liski and coauthors simulated the soil C of Finland's forests
forward from 1922 to 2004, over nearly the same period, with the
same model family, and from litter inputs derived from the same National Forest
Inventory. But Liski and coauthors had no repeated measurement of
the same soil, from which trajectory we could constrain an initial state, so they assumed one,
calculating the pools at 1922 ``by assuming a steady state with mean litter input between
1922 and 1936 and mean temperature between 1901 and 1930''. The present study keeps their
forward simulation, but replaces that one assumption with an estimate that draws information from the measured SOC trajectory.

\subsection{A managed forest landscape is not an equilibrium landscape}
A forest whose standing biomass has changed within the time the soil needs to adjust to a
new litter input leaves a soil that is not in balance with that input, and this is a rather
common occurrence. Finland is a well-documented case
of one. The Finnish forest of the early twentieth century had been heavily used:
slash-and-burn cultivation had passed over most of the productive forest land of eastern
Finland, grazing in the forest suppressed regeneration, and tar burning and the selective
removal of the best timber had drawn down the accessible stands. By the first national
forest inventory old trees were effectively absent from southern Finland outside the least
accessible areas, and the volume of wood in Finnish forests had reached its historical low
point \citep{Aakala2023}. The recovery that followed was not steady. Logging fell during the
war years and the growing stock rose in the 1940s, but the use of wood for reconstruction and
export reduced it again through the 1950s and 1960s, and the strong increase in logging before
the fourth inventory is visible in that inventory as a slowdown of the increment
\citep{Korhonen2021}. The pressure had three sources: the war reparations paid to the Soviet
Union until 1952, about four per cent of the national product each year, delivered as
industrial goods but financed by an economy whose exports came from the forest sector
\citep{Mitrunen2024}; the loss in 1944 of a ceded territory that held more than a tenth of the
country's forests and a fifth of its pulp capacity \citep{Kochetkova2020}; and the resettlement
of the evacuated population on land taken from the forest. Only after the late 1960s did the
growing stock rise without interruption, from roughly $1.4$ to $2.6$ billion cubic metres over
the century as a whole, an increase of $84$ per cent \citep{Korhonen2024}
(Figure~\ref{fig:history}).
% LORENZO note (2026-09-16) on the war-reparation pressure: answered 2026-09-21. Sources found
% online: Korhonen et al. 2021 (Silva Fennica 10662, the NFI's own record of the 1950s-60s
% drawdown), Mitrunen 2024 (QJE; reparations 4% of GDP, mostly metal goods -> "financed", not
% "paid in timber"), Kochetkova 2020 (EHNE; ceded Karelia 12% of forests, 20% of cellulose
% capacity). Henttonen et al. 2020 checked and dropped: it skips the war decades.
% Details: literature/04_finnish_management_history/README.md, "Update 2026-09-21".

\hfig{F10b_forest_history}{0.92}{history}{\textbf{The growing stock of Finnish forests,
1921--2023.} Growing stock rose by $84$ per cent over the century shown, and most of the rise
occurred after the late 1960s. The litter falling to the soil is tied to the standing biomass, so
the litter input that the soil has integrated has been rising throughout the period that the models
are asked to reproduce. Data from \citet{Korhonen2024}.}

But soil C responds slowly to a change in litter input. The pools that turn over within
years adjust within a few decades, whereas the slow pool that holds most of the stock
approaches its new level only over centuries, so that a soil disturbed several centuries
earlier can still be accumulating \citep{Wutzler2007}. The move from a C-poor soil under a
low litter input to a richer soil under a higher one is therefore measured in decades to
centuries, not years. That time is much longer than the mean time a unit of carbon spends in
the soil, because the carbon leaving the soil comes mostly from the fast pools while the
carbon stored in it sits mostly in the slow one; the two timescales are distinct properties
of the same model \citep{Sierra2018}.
% LORENZO note (2026-09-16) "this statement should be referenced. How much time?": answered
% 2026-09-21 with the pool-specific timescale (decades for the fast pools, centuries for the
% slow pool) on Wutzler & Reichstein 2007; Liski et al. 2006 in the next paragraph is the
% Finnish instance ("centuries later"). No own computation of the slowest-mode e-folding time
% (decided 2026-09-21); literature/03_equilibrium_theory/README.md sec. C describes it if wanted.

The consequence for the starting state of a soil model was stated in this literature two decades ago. 

\citet{Peltoniemi2004}, working from soil C stocks measured onlyonce on Finnish stands of different ages, showed that such a chronosequence cannot separate the two processes that act on the soil at the same time: the change in soil C over the course of a rotation, driven by management, and the slower recovery of the whole landscape from its historical depletion. 

\citet{Liski2006} approached the long-term trend by simulation, and stated both its direction and its shape: Finnish forest soils would go on accumulating C even if the litter input were held at its 2004 level, and, centuries later, the stocks would stabilise at a level about $38$ per cent above that of 1922. An accumulation that stabilises is an accumulation that slows dowhn, and the slowing is the part of that statement which has received the least attention since.

Finland now has three soil C campaigns on the same plots, spanning roughly four decades, which measures the trend directly, and that trajectgory is the information that can make the initial state something to be estimated rather than assumed.

\subsection{What a soil started in balance cannot do}
% Rewritten 2026-09-16 after the discussion with Lorenzo. Three claims, in order:
% (1) a stock fixes only the product J x MTT (Olson) -- equifinality, named, Beven;
% (2) a trajectory narrows that only through its CURVATURE (the slowing), which
%     identifies the turnover of the pools that respond within the window, and only
%     as well as the second-interval rate is measured -- so it narrows, weakly;
% (3) a balanced start does NOT lose the ability to slow down (the old sentence
%     claimed that, and it was false: it responds to any input change after the start);
%     what it loses is the MEMORY -- the deficit already present on the start date.
%     With anchored rates that deficit is not free (the product is set to the observed
%     stock; a soil still catching up implies a larger one -- ours 28-49% larger);
%     with free rates the history is stored in the rates, which the projection runs on
%     (Wutzler & Reichstein 2007). Luo 2017 dropped: it did not support the old claim.
% What is NOT said here, because it has not been tested: whether the ridge is wider
% under a balanced start, and whether the inventory's convention with the actual litter
% series reaches the observed rate/level ratio. Both are the equilibrium-init
% counterfactual (NEXT_RUN_equilibrium_init.md).

Two quantities govern the C stock of a decomposition model: the litter input the soil
receives, and the mean transit time of C in the soil. A measured stock does not separate
the two, because the same stock is reached by a large litter input decomposed quickly and
by a small litter input decomposed slowly. A soil in balance holds a stock equal to the
product of the two quantities, $C^{*} = J/k$ for a single pool \citep{Olson1963}, so the
product is all that a stock can fix. The same observation is reproduced by many
combinations of input and kinetics, which is the situation \citet{Beven2006} named
equifinality, and holds for every soil C model and for the transient phase
of a simulation as much as for its steady state.
A trajectory narrows that set, but only through its shape. For a single pool the ratio
of the rate of change in one interval to the rate in the interval before depends on the
decomposition rate alone, and not on the litter input nor on how far from balance the
soil started, so the variation of such rate over time carries additional information.

A soil started in balance is not unable to accumulate or to slow down, it still responds to
 the litter input and the climate after the start. What it cannot carry is a
disequilibrium that already existed on the start date. With decomposition rates held at their
published values, there are limits to the shape of the trajectories the model can represent and the model would not be able to fully represent a trend that originates also in an initial disequilibriu,. But if the rates are freed so that a balanced start fits the
observed trajectory anyway, the missing history would be absorbed into the rates
\citep{Wutzler2007}, and this might result in biased future extrapolations.

\subsection{Objectives}

This study revises two assumptions about the starting state, and the two work together. 

First,
the state is estimated rather than assumed: the litter input at which the soil is taken
to be in balance is a free parameter, and the measured trajectory is what informs it,
because a stock alone cannot. 

Second, the date of that balance is pushed back to 1917,
as far as the forest record allows, so that whatever is wrong with the assumption of
balance in 1917 (and something will likely be, since a real steady state is something achieved only in unmanaged, primary forests) has sixty-eight years of simulated history to dissipate that error while being influenced by the transient simulation before the first
measurement. 

A starting
state assumed in balance just a decade before the first measurement has neither protection.
This makes possible to answer the questions below.

I) Once the
starting state is estimated from the measured trajectory, where does that trajectory
place the models on the ridge, in parameter space (and defined by similarly good fit to the data), between litter input and transit time, relative to
their published parameterisations? And how far does it narrow that exchange?

II) Does the position on such ridge matter beyond the observed
window? Do two calibrations that reproduce the same measurements from different
positions extrapolate differently, and for which model structures? 

III) What part of
the variation between individual plots (due to local conditions) remains unexplained, and what 
information would the model need to explain it?

% =============================================================================
% ARCHIVE. The Introduction as it stood before the 2026-09-07 restructure,
% verbatim, with your original notes in place. Delete once the new version has
% been read and accepted.
% =============================================================================
% \section{Introduction}
%
% \subsection{A soil carbon sink that appears to be levelling off}
%
% Finland has measured the carbon stock of the same forest soil plots three times
% over four decades. 
% %this needs more qualification about the national inventory efforts. But I am not convinced about it being the opening, even if it is indeed the core of the problem, it is a bit too method specific. I would state the problem more generically in the opening phrase, going then specific about the samplings.
%
% On the set of plots that were measured in all three campaigns,
% the mean soil carbon stock rose from $64.5$ tonnes of carbon per hectare in the
% first campaign to $71.5$ in the second and $73.6$ in the third. The stock
% therefore gained $7.0$ tonnes of carbon per hectare over the first interval and
% $2.1$ tonnes over the second interval. The soil accumulated carbon over both
% intervals, and the accumulation was slower over the second interval than over the
% first. 
%
% \hfig{F3_mean_SOCalt}{0.86}{observed}{\textbf{The measured soil carbon stock of Finnish forest plots in three campaigns.} Each point is the mean stock across the plots measured in all three campaigns ($n=310$), and the bars show the $95\%$ confidence interval of that mean. The first campaign is drawn at the mean year in which its plots were actually sampled, and the bar low in the panel shows that the sampling of the first campaign ran from 1986 to 1995. Each segment is labelled with the rate of change over that interval and with the $95\%$ confidence interval of the rate, calculated from the paired difference between campaigns on the same plots. The stock rose over both intervals, and the rise over the second interval was smaller than the rise over the first.}
%
%
% Expressed as a rate, the mean stock rose by $0.259 \pm 0.040$ tonnes of carbon per
% hectare per year between the first and the last campaign.
% % connect these two with a but
%  Over the last interval
% considered on its own, the mean stock rose by $0.117 \pm 0.066$ tonnes of carbon
% per hectare per year, and the confidence interval of the rate over the last
% interval includes zero. 
% %last phrase about the CI including zero seems to mean nothing, also becose CI was not even mentioned so far. Maybe remove? Or is it functional there to something?
% Every rate reported in this paper is calculated on the
% same basis, which we state once here and repeat wherever a rate appears: the set
% of plots measured in all three campaigns ($n = 310$), each plot weighted equally,
% the whole soil profile, and the true year in which each plot was sampled.
% % THis above is again a bit too much M&M
%
%
% %Here below maybe intro phrase about the difference between state and trajectories, and the latter is what matters for us (and why)
% A soil that is still accumulating carbon and a soil that has stopped accumulating
% carbon hold the same amount of carbon on the day they are measured, but the two
% soils have different futures. For a greenhouse gas inventory, which must project
% the stock forward, 
% %not necessarily that, even if it is also important for policies. The key is that the national budget is calculated on the sink/source, which is determined by the rates (trajectories not state).
% the quantity that matters is therefore not the stock itself but
% whether the accumulation of the stock is slowing down.
%
% \subsection{The soil term in the national inventory is produced by a model}
%
% Finland reports the carbon balance of its forest mineral soils to the United
% Nations Framework Convention on Climate Change using a modelling method rather
% than a measurement. 
% %... as many other countries. Here short review, possibly exhaustive, of national budgeting methods.
%
% Forest inventory data and biomass models supply an estimate of
% the litter falling to the soil, and
% %, in the case of Finland,
%  the decomposition model Yasso07 converts that
% litter estimate into an annual change in the carbon stock of dead organic matter
% and soil organic matter combined, to a depth of one metre, calculated separately
% for southern and for northern Finland
% \citep{StatFin2024NID}. The method follows \citet{Liski2006}. No soil carbon flux
% is measured at any point in the reporting chain.
%
% The soil term is large enough to matter and uncertain enough to matter. 
% %this phrase sounds cool but it is a bit doubtful, particularly on the meaning of "uncertain enough to matter". Might be fine as concept, maybe, but we need to sharpen it if we keep it
% In 2022
% the mineral soils of Finnish forest land were reported as a sink of $4.8$ million
% tonnes of carbon dioxide, alongside a living-biomass sink of $12.6$ million
% tonnes, while the forest land sink as a whole had fallen by $84$ per cent since
% 1990 and the land use sector had become a net source of $4.4$ million tonnes of
% carbon dioxide equivalent \citep{StatFin2024NID}. 
% %Ok, now here we get to the core of the problem. Maybe the intro is a bit too long and not sharply focused, this is the tension
%
% The reported uncertainty of the
% mineral soil term is $31.5$ per cent, against $16.3$ per cent for the tree biomass
% term it accompanies \citep{LehtonenHeikkinen2015}. When the national balance is
% close to zero
% %meaning what, the modeled balance?
% , a model-derived term of that size, carrying that uncertainty, can
% determine the sign of the reported balance.
% %this is a bit self evident, like the number determines that same number
%
% The reporting chain also 
% %also in respect to what? Are here enumerating some of the possible pitfalls? THis should be done in a more logically consistent way. And each of those pitfalls (related to some assumptions) has practical implications that are important, and we explain which and why
% fixes the state of the soil at the start of the
% simulation, and the chain states the assumption it makes when it does so. The
% initial state is obtained by running the model forward under an unchanging litter
% input until the carbon stock stops changing, using the average litter input of the
% sixth National Forest Inventory (1971--1976) and the average weather of 1960--1990.
% The justification given is that ``approximately ten years of simulation since
% spin-up is enough to cancel out the effect of the spin-up level''
% \citep[\S6.4, citing][]{StatFin2024NID,Peltoniemi2006}. The present paper tests
% that justification over a window of thirty-five years, using an initial state that
% is estimated from the soil measurements rather than assumed.
%
% \subsection{Starting a soil model at steady state is the standard practice}
%
% Starting the simulation from a soil that is already in balance with its litter
% input
% %steady state model initialization. THis would not be necessarily jargon, it is quite used. "Running a model into a steady state" maybe is a bit weird construction (would not cite that one, seems we are making fun of it), but "running it" is the weird word here, not steady state. Steady state assumption in SOC models initialization is standard way of calling it.
%  is the standard practice across the model families used for forest soil
% reporting. In the Finnish inventory, the models Yasso07 and ROMULv 
% %from where now ROMULv comes from? First time it is mentioned? Is this because we are listing some examples of steadystate initialisations, arguing that it is common?
% were ``run into
% a steady state'', on the stated premise that ``if we assume that this average level
% of inputs and climate has remained steady over centuries, then our soils should
% approach steady-state conditions'' \citep{Lehtonen2016}. In applications of the
% model RothC, the same equilibrium run is used additionally to work backwards from
% the measured stock to the annual carbon input \citep{Contreras2026}. Global land
% models are initialised in the same way \citep{Luo2016}. The Canadian forest carbon
% model CBM-CFS3 is the instructive exception, because CBM-CFS3 initialises the soil
% by simulating repeated cycles of growth and disturbance and therefore reports an
% initial state that is deliberately not in balance with the current litter input
% \citep{Kurz2009}.
% %Is this so radically different? If the cycles are oscillating around a mean, it is just a bit less approximate but conceptually does not seem that far
%
% The closest predecessor of the present study is \citet{Liski2006}. Liski and
% colleagues simulated the soil carbon of Finland's forests forward from 1922 to
% 2004, driven by litter inputs derived from the national forest inventory, using
% the same model family and the same country as the present study. 
% %and here instead it is indeed a different conceptual approach, pretty much the same we propose here
% Liski and
% colleagues did not, however, estimate the state of the soil at the start of the
% simulation. In their own words, ``the soil and litter carbon pools at the
% beginning of the study period were calculated by assuming a steady state with mean
% litter input between 1922 and 1936 and mean temperature between 1901 and 1930.
% Starting from this steady state in 1922, the model was run using annually varying
% values of litter input and temperature'' \citep{Liski2006}. The present study
% keeps the forward simulation and estimates the starting state instead of assuming
% it. One assumption separates the two studies, and repeated measurements of the
% same soil are what make it possible to replace that assumption with an estimate.
%
% Starting a soil model at steady state has been criticised before, and the
% criticism is not what this paper contributes.
% % maybe the last part of the phrase is too open about our intentions, even if this information (that we are not aiming at that specifically) could be part of the intro
%  \citet{Wutzler2007}, working with the
% Yasso model, showed that if a soil is in fact still accumulating carbon, then the
% stock it would eventually reach becomes very large, 
% % this above makes little sense as it is. A soil could be accumulating C and then the model would saturate, even pretty close to it. Wutzler probably meant something else. Does it involve some fitting, in certain conditions? If we fit on measurements of a soil accumulating, breaking the steady state assumption clearly?
% and that fitting model
% parameters to the present stock under the assumption of balance therefore
% overestimates how fast the slowest carbon pool decomposes. \citet{Carvalhais2008}
% showed that the same assumption biases the parameters retrieved by fitting a model
% to data, and not only the initial state. Later studies have measured what the
% choice of initial state costs in predictive error
% \citep{Lee2020,Kanari2022}. What has been missing is not the criticism but the
% test.
% %Ok, so we aim at testing such criticism. Is it too ambitious? But maybe it is fine, I would avoid in this case to mention that we do not want to criticise, it seems "don't think to the elephant" kind of thing
%  Wherever the assumption has been relaxed, the replacement initial state has
% been imposed --- from an assumed disturbance history, from a laboratory proxy, or
% as a prescribed correction --- rather than estimated from an observed change in
% the carbon stock itself.
% %Really? And, for example, Liski before?
%
% \subsection{What starting at steady state assumes, and what follows from it}
% In all of the models used here, a carbon pool loses carbon in proportion to the
% amount of carbon the pool holds, and in proportion to nothing else. Under a
% constant litter input and a constant climate, a system of pools of that kind
% drifts towards a single carbon stock and then stays there. The stock it reaches is
% the litter input multiplied by the mean time a carbon atom spends in the system,
% which for a single pool is the familiar expression $C^{*} = J/k$
% \citep{Olson1963}. Starting a simulation at that stock is not an assumption about
% the soil. Starting a simulation at that stock is an assumption about the history of
% the soil: that the litter input and the climate have been effectively unchanged
% for long enough that even the slowest carbon pool has caught up with them, and that
% no disturbance has interrupted the soil in the meantime.
% %this theory is good, but maybe a bit too disperse, can be denser. Also maybe it's not the right place, I would introduce what it thsi assumption is earlier, when explaining how many models do it and how it is standard practices. Consequences instead deserve its own paragrpah
%
% Two consequences follow from the assumption, and the present study turns on both
% consequences. 
% %"Turns on" unclear, more plain language
% The first consequence is that the condition 
% %you mean assuming steady state, with "condition of balance"?
% of balance fixes only
% the product of the litter input and the mean residence time. 
% %I donät get what this means, it is about calibration? That at steady state there is a narrow edge of combinations in the parameter space that can achieve that state (between inputs and kinetics)?
% A given stock can
% therefore be reached by a soil that receives a large litter input and decomposes it
% quickly, or by a soil that receives a small litter input and decomposes it slowly,
% and a measurement of the stock alone cannot distinguish between the two soils.
% %yes, seems setting the equifinality problem, our ridge, that indeed exisist in all these models
%  The
% trade-off 
% %is it really a trade-off, that ridge?
% between litter input and residence time is a property of the model
% equations, and the trade-off is present before any data are fitted. The second
% consequence is that a soil started in balance has no accumulation left to perform.
% A soil started in balance can follow a target that moves as the climate and the
% litter input change, but a soil started in balance cannot express a soil that is
% still catching up with its litter input, and therefore cannot produce an
% accumulation that slows down \citep{Luo2017}.
%
% \subsection{Finnish forests are a managed landscape and not an equilibrium one}
% %Ok, here follows the rationala about the criticism on the steady state assumption for the specific case of Finland
%
% The Finnish forest of the early twentieth century had been heavily used.
% Slash-and-burn cultivation had passed over most of the productive forest land of
% eastern Finland, grazing in the forest suppressed regeneration, and tar burning and
% the selective removal of the best timber had drawn down the accessible stands. By
% the first national forest inventory, old trees were effectively absent from
% southern Finland outside the least accessible areas, and the volume of wood in
% Finnish forests had reached its historical low point \citep{Aakala2023}. A century
% of managed regrowth followed. The growing stock of Finnish forests rose from
% roughly $1.4$ to $2.6$ billion cubic metres between the first and the thirteenth
% inventory, an increase of $84$ per cent, and most of the increase occurred after
% the late 1960s \citep{Korhonen2024}. The soil modelled in this study has integrated
% the whole of that century (Figure~\ref{fig:history}).
%
% \hfig{F10b_forest_history}{0.92}{history}{\textbf{The growing stock of Finnish forests,
% 1921--2023.} Growing stock rose by $84$ per cent over the century shown, and most of the rise
% occurred after the late 1960s. The litter falling to the soil is tied to the standing biomass, so
% the litter input that the soil has integrated has been rising throughout the period that the models
% are asked to reproduce. Data from \citet{Korhonen2024}.}
%
%
% The litter that falls to the soil is tied to the standing biomass of the forest, so
% the litter input that a soil model must integrate has been rising throughout the
% period the model is asked to reproduce. Whether the soil has kept pace with the
% rising litter input is a separate question, and the available evidence indicates
% that the soil has not kept pace. Old unmanaged stands in southern and eastern
% Finland hold about $40$ per cent more carbon in the mineral soil than managed
% stands of comparable site types, while the carbon stocks of the unmanaged stands
% themselves did not change measurably over a twenty-year resampling interval
% \citep{Kumpu2025}. A managed landscape that holds less carbon than its unmanaged
% counterparts, under a litter input that is still rising, is a landscape whose soil
% is not in balance with its litter input.
%
% The consequence for the initial state of a soil model was stated in this literature
% two decades ago. \citet{Peltoniemi2004} attributed part of the soil carbon of
% Finnish forest plots to a history of slash-and-burn cultivation that ended in the
% early twentieth century, and concluded that ``the long-term trend in carbon
% accumulation cannot be distinguished from the measured data unless the measurements
% cover at least two rotations''. \citet{Liski2006} reported that Finnish forest
% soils would continue to accumulate carbon even if the litter input were held at its
% 2004 level, and that the stocks would settle ``centuries later \dots\ at a $38$ per
% cent higher level than in 1922''. Finland now has three soil carbon campaigns
% spanning roughly four decades, which is what makes the initial state of the soil
% estimable from the measurements.
%
% \subsection{Objectives}
%
% The convention of starting at steady state is the reason the ability of these
% models to follow a national soil carbon trajectory has not been tested. A soil
% started in balance reproduces a stock but has no trajectory that could be shown to
% be wrong. Replacing the assumption is therefore not a methodological preliminary to
% the results of this paper. Replacing the assumption is what allows the results to
% exist at all.
%
% We ask three questions, and we weight them in the order in which they are listed.
% First, does a transient and calibrated initial state reproduce the observed
% national soil carbon trajectory, and what does reproducing the trajectory require
% of the models? Second, does it matter, for projection beyond the observed window,
% where a calibration lands on the trade-off between litter input and residence time?
% Third, what part of the variation between individual plots can none of the models
% see, and what would be needed to see it?
%
% =============================================================================

% =============================================================================
\section{Materials and methods}

% Carried over from HIKET_main_manuscript.tex and checked against the code on
% 2026-09-04. Every setting stated below was read from calib_config.R, the SLURM
% submission scripts of the production run, correlated_likelihood.R, the prior
% specification files, or the input bundles themselves.

\subsection{Study plots and soil carbon observations}

The study is based on data from permanent forest plots of the Finnish National Forest Inventory for
which soil carbon has been measured in three campaigns: the eighth national forest
inventory (VMI8), the BioSoil campaign of 2006, and the Komeetta campaign of 2024.
After the exclusions described below, $456$ plots are available for calibration,
carrying $1205$ plot-year observations of soil carbon in total. Of the $456$
plots, $365$ plots ($80$ per cent, carrying $959$ observations) were used to fit
the models, and $91$ plots ($20$ per cent, carrying $246$ observations) were held
out and used only for evaluation. Plots were assigned to the two groups at random,
and a plot was assigned as a whole, so that no plot contributes observations to
both groups.

The first campaign is not a single year. Its plots were sampled region by region
between 1986 and 1995: 69 plots in 1986, 76 in 1987, 103 in 1988, 54 in 1989 and
72 in 1995. Averaged over the plots, the sampling year of the first campaign is 1989,
and that is the year at which the campaign is drawn whenever it is shown as a single
point. In the calibration, each observation is compared with the model in the year in
which that plot was actually sampled.

On the plots measured in all three campaigns ($310$ plots, the set on which every
observed rate of change in this paper is computed), the mean stock of the groups sampled
in 1986, 1987, 1988, 1989 and 1995 differs by $18.3$ tonnes of carbon per hectare at
first visit, which might reflect a trend within that same campaign.
% ROUND 5, 2026-09-23: "the same groups" named explicitly (note 1044); the
% closing sentence, which repeated the previous paragraph, deleted.
The five groups of plots defined by their first-campaign sampling year (1986, 1987,
1988, 1989 and 1995) differ by $37.2$ tonnes in 2006 and by $34.2$ tonnes in 2024,
years in which every plot of every group was measured in the same campaign. The
difference between the five groups is therefore a property of which plots each group
contains, and not of the year in which it was sampled, because the difference is
larger when the sampling year is held constant.

\subsection{Soil carbon stocks on a single homogenised basis}

% CARRIED from the previous manuscript; the canonical text of this argument lives
% in the METHODS AND MATERIALS block at the end of Data/Data_work.R.
% ROUND 2, 2026-09-11, in answer to your notes:
%   * "which workbook" -> NOT named, and no citation. (Lorenzo, 2026-09-11: an
%     unpublished internal compilation is not a reference; citing it would dress our
%     own data up as a source.) The text now says what the data ARE -- the layer-wise
%     measurements of the three campaigns, put on one set of conventions here -- and
%     the provenance and acknowledgement belong in the data availability statement.
%     A \gap marker further down carries that.
%   * stoniness, "I fear we don't" -> WE DO, but not by correcting it ourselves: the
%     stocks in that workbook are already corrected for coarse fragments and bulk
%     density. The paragraph now says what we take, not what an earlier assembly
%     omitted.
%   * "exactly" -> I re-ran Data/SOC_homogeneized/build_soc_homogenized.R. On the
%     official calculation's own exclusions (n=446) it returns 59055 and 61047 kg/ha
%     against targets of 59055 and 61047, i.e. agreement to the nearest kg/ha. The
%     text now states the check instead of asserting the adverb.
%   * "which earlier assembly" -> the comparison to an earlier assembly is GONE.
%   * layer standardisation and sampling years -> rewritten as standardisation and
%     as a choice, not as corrections, in your words.

Since the three campaigns were measured and processed under protocols
that differed in ways that matter for a comparison between them, all three campaigns
were rebuilt on a uniform basis, starting from the layer-by-layer carbon
stocks measured in each of the three campaigns and applying one set of conventions to
all three. Those stocks are corrected for the volume occupied by stones and for the
bulk density of the fine earth.
% ROUND 5, 2026-09-23 (note 1081): methods from Lehtonen et al. 2016 (GMD) and
% Kramarenko 2012, sect. 2.3-2.4; the 1985 stocks use the 2006 stone volumes.
The volume of stones was measured in 2006 on every plot by the rod-penetration method
of \citet{Viro1952}, in which a steel rod is driven into the mineral soil at the
sampling points and the mean depth it reaches is converted into a stone volume, and
the same stone volume is applied to the first campaign \citep{Lehtonen2016,Kramarenko}.
The bulk density of the mineral layers was predicted from the equations of
\citet{TamminenStarr1994} for all layers of the first campaign and for most layers of
2006, and measured directly for part of the 2006 samples from the upper $20$ cm
\citep{Kramarenko}. \gap{Confirm that the 2024 campaign uses the same stone volumes
and the same bulk density equations.}
The coarse-fragment correction matters
because the coarse-fragment content of these plots has a median of $43$ per cent by
volume, so a mineral stock computed without it is too large by a factor of about $1.6$.

The quantity used as the calibration target is the whole soil profile: the measured
organic layer, plus the measured mineral soil to $40$ cm, plus a modelled deep
tail, integrated only as far down as the auger actually reached on that plot and in
no case deeper than one metre. The depth to which the profile is integrated is
therefore a per-plot quantity and not a fixed one metre, since using a fixed depth of, for example, one metre
would invent soil carbon below bedrock on the $38$ plots where augering stopped at
$10$, $20$ or $40$ cm. The modelled deep tail was validated against the measured
$40$--$80$ cm layer of the BioSoil campaign, which was held out of the fit: the
median ratio of predicted to observed carbon in that layer is $0.96$, with a bias
of $-2.5$ tonnes of carbon per hectare.

The rebuilt target was checked against the national figures published for the two
campaigns for which such figures exist. Applying the same plot exclusions as the
official calculation, and integrating over the same layers, the rebuilt target returns
weighted national means of $59.06$ and $61.05$ tonnes of carbon per hectare for 2006
and 2024 against official values of $59.06$ and $61.05$, that is, agreement to the
nearest $0.01$ tonnes per hectare on $446$ plots.

Two further steps were taken with the first campaign, both of them established from the
structure of the measurements themselves rather than from the fit. First, the layers of the
three campaigns were standardised, because they were recorded differently: the organic
layer of the first campaign contains the fermentation and humus horizons only, whereas
the organic layer of the two later campaigns also contains the litter and moss layer.
Each plot's own litter and moss layer from 2006 was therefore added to its
first-campaign organic layer, so that the same compartment is compared across the three
campaigns. Second, the first campaign is treated as what it is, a campaign whose
fieldwork ran over several years rather than a campaign carried out in a single year.
The year in which each plot was actually sampled is on record, and those years,
rather than one nominal year for the whole campaign, are the years at which the models
are compared with the observations.

\gap{Data availability and acknowledgement: the soil carbon measurements are not a
published dataset, so the statement has to say where they come from, who compiled them
and what of the compiled form can be deposited. Settle this with the people who hold
them before submission.}

On this final basis the median whole-profile stock is $62.0$ tonnes of carbon per
hectare for the first campaign, $66.4$ for 2006 and $67.2$ for 2024, and the
observations vary between plots with a coefficient of variation of $0.44$ to $0.53$
depending on the campaign.

\subsection{Litter input}

% ROUND 5, 2026-09-23 (notes 1130, 1136): product named and its method stated from
% the Tupek et al. manuscript in the Zenodo deposit, sect. 2.7-2.8. CORRECTION: the
% product has NO natural-mortality term (turnover + harvest residue only; a tree that
% dies registers only through a basal-area drop, treated as a harvest with the stem
% removed). The old sentence "includes ... natural mortality" was wrong.
Litter input is taken from the annual reconstruction of tree biomass and litter input
for the permanent plots of the Finnish National Forest Inventory by
\citet{TupekLitter}, published as an open dataset. In that reconstruction, the tree
biomass measured by the inventory is interpolated to every year with the help of
satellite-derived basal area, and litter is computed as the biomass of each tree
component multiplied by a turnover rate for that component, plus the branches,
foliage, stumps and roots left on site when a drop in basal area identifies a
harvest; the full method is given in \citet{TupekLitter}. The product supplies, for
each plot and year, the carbon in litter from trees, separated into non-woody litter
(foliage and fine roots), fine woody litter (branches and coarse roots) and coarse
woody litter (stumps), and each of the three classes is further divided into the four
chemical fractions used by the Yasso models, using proportions specific to each tree
species and component. The product does not contain a separate term for trees that
die without being harvested, and it does not include the understorey vegetation,
because the product is a tree litter product by design. Across the study plots the mean litter input carried
by the product is $\bar{J} = 2.511$ tonnes of carbon per hectare per year.

The first year of the product is treated separately. The median litter input of 1985
is $0.060$ tonnes of carbon per hectare per year against $1.398$ in 1986, a step that
the series does not repeat in any later year. Because 1985 is both the first year of
the simulation and the year in which the first soil observations fall, the value for
1985 was reconstructed for each plot and each chemical fraction by fitting a linear
trend over 1986--1990 and extrapolating one year backwards. The reconstruction places
1985 just below 1986, which is the direction the growing-stock record requires.



\subsection{Climate}
Climate forcing comes from the gridded daily weather observations of the Finnish
Meteorological Institute, interpolated from the station network onto a $10$ by $10$ km
grid by kriging with external drift \citep{Aalto2016} and extended by the Institute to
the present. The grids were obtained through the internal weather database of the Natural
Resources Institute Finland, which mirrors them from the Finnish Meteorological
Institute. From the grid cell nearest to each plot we took the daily mean, minimum and
maximum temperature, precipitation and global radiation for 1961--2025, and each model
takes from those series the variables it requires. Because the response of
decomposition to climate is not a straight line, the average of the response over the
years is not the response evaluated at the average climate. Averaging the climate first
and evaluating the response afterwards biases the climate modifier by $5$ to $15$ per
cent in boreal conditions. The climate modifier is therefore computed for each year and
averaged afterwards, in all six models.

\subsection{The six models}\label{sec:models}

Six models are used, and their roles are different by class. 
Yasso is the model this
study is focused on, because Yasso is the model the Finnish greenhouse gas
inventory runs. All three published versions of Yasso are included: Yasso07,
Yasso15 and Yasso20. 
Alongside the three versions of Yasso we included three simple
models with one, two and three carbon pools, denoted SP1, TP2 and TP3. The three
simple models are there as a benchmark, since their behaviour is transparent enough so that any result obtained with Yasso can be checked against a
structure simple enough to reason about quickly.

Due to their different structures and designs, the number of free parameters differs greatly across the six models, from six free
parameters in SP1 to twenty-six in Yasso15 and Yasso20. The free parameters of the
three versions of Yasso are the twelve fractions describing transfers between
pools, the climate parameters, the parameters describing the effect of the size of
woody litter, and the two auxiliary parameters described in Section~\ref{sec:inputunc}. 

The decomposition
rates of Yasso are not fitted, and are held at their published values.
Holding the rates fixed is a deliberate choice to keep the model anchored to its original definitions.
Decomposition rates and transfer fractions do not act independently on the carbon
stock but are always interacting: a slower rate and a larger transfer into a slow pool produce much the same
trajectory, so fitting both together improves usually little and makes the result harder to read.
% ROUND 5, 2026-09-23 (notes 1195, 1198, 1204): the literature digression, framed
% descriptively (Yasso's own transfers, no verdict), reconciling compartments and
% continuous quality through fractionation-as-discretization. The Yasso07 numbers are
% the published p_NA = 0.970 and p_AW = 0.987 (Prior_specs/Yasso07_priors.R, from
% Tuomi et al. 2009); direction checked in yasso07.f90, A(dest, source).
Of the two, the rates are the more fundamental to the definition of each model,
because the rate of a pool is what distinguishes that pool from the others. Whether
the pools of a soil carbon model can also be identified with measurable fractions of
organic matter is a long-standing question in soil carbon modelling, and the answer
has shaped how Yasso was built. The first version of Yasso assumed that litter
consists of groups of organic compounds, each group decomposing at its own rate
regardless of the litter it came from, and defined its decomposition compartments
accordingly as extractives, celluloses and lignin-like compounds, followed by two
compartments of humus \citep{Liski2005}. Yasso07 kept the same principle and defined
its four litter pools by fractions that can be measured on litter by sequential
extraction: the fraction soluble in acid, the fraction soluble in water, the fraction
soluble in ethanol, and the insoluble remainder \citep{Tuomi2009}. Defining the pools
by the chemistry of the litter is what makes Yasso operational, because the chemistry
of litter can be measured and supplied to the model directly. The same aim, of
matching the pools of a model with fractions that can be isolated in the laboratory,
has been pursued for other models, for example through fractionation schemes designed
to correspond to the pools of RothC \citep{Zimmermann2007}, and whether fractions
isolated in the laboratory correspond to the functional pools of models has been
examined at length \citep{vonLutzow2007}.

Once carbon has entered the soil, the parameterisation of Yasso moves carbon between
the four chemical fractions in every direction. In the published parameterisation of
Yasso07, for example, $97$ per cent of the carbon decomposed from the insoluble
fraction is passed to the acid-soluble fraction, and $99$ per cent of the carbon
decomposed from the acid-soluble fraction is passed to the water-soluble fraction
\citep{Tuomi2009}, which represents the transformation of organic matter by
decomposers rather than the decay of a fixed set of compounds. The pools therefore
enter the model defined by the chemistry of the litter, and within the soil they act,
as in any other compartmental soil carbon model, as compartments defined by their
rates of turnover.

The view of pools as compartments of turnover and the view of soil organic matter as
a continuum of quality are less opposed than they appear. A chemical fractionation is
itself a discretization: it divides a continuum into a small number of classes, the
discretization is made in the laboratory rather than in the model, and where the
boundaries between classes fall is set by the method chosen, so that the fractions are
defined by the procedure that isolates them \citep{vonLutzow2007}. The material being
divided is not a set of discrete compounds. Soil organic matter is the product of many
processes of decomposition, microbial transformation and recombination, much of that
organic matter has no defined molecular structure, and the chemical diversity of soil
organic matter is large enough that the diversity can be treated as a continuum
\citep{LehmannKleber2015}.

The theory of continuous quality starts from that continuum: \citet{BosattaAgren1991}
describe soil organic matter by a distribution of quality, in which the quality of a
unit of carbon determines how fast that unit is decomposed, and decomposition in turn
moves the carbon along the continuum \citep{AgrenBosatta1998}. The same idea was
formalised independently in the geochemistry of sediments, where models with a few
discrete fractions of different reactivity are treated as discretizations of a
continuous distribution of reactivity \citep{BoudreauRuddick1991}. On this view, a
pool of a compartmental model and a fraction of a laboratory procedure are two
discretizations of the same continuum, one made in the model and one made in the
laboratory, and in the model the rate of decomposition is what places each pool on the
continuum. Changing a rate therefore redefines the pool, whereas the transfer
fractions describe how carbon moves between pools that have already been defined. We
therefore keep the rates as the published parameterisation set them and calibrate the
transfer fractions. In all three versions of Yasso the decomposition rates, together
with the fraction of decomposed carbon that enters the humus pool, are fixed at their
published values: they are constants of the model and not parameters of the
calibration, and no prior is placed on them.

\subsection{The uncertainty of the litter input}\label{sec:inputunc}

% ROUND 5, 2026-09-23 (notes 1260, 1284): the two sigmas moved ahead of the transient
% initialisation and presented as the treatment of input uncertainty, extended to the
% start. Framed as INVERSE vs the FORWARD propagation of Peltoniemi 2006 / Lehtonen &
% Heikkinen 2015 (verified from the L&H PDF); no "rarely done" claim. sigma_init
% centre 0.90 = 0.789^0.43 (Prior_specs); Liski 2006 per-ha J1922/J1985 = 0.861 with
% the 1965-1980 area timing (0.852 vs the 1985-89 mean): a consistency check on the
% same NFI record, not an independent source, and not used in the calibration.
The litter input is among the least certain of the quantities that drive a soil
carbon model. Litter is not measured on the plots but derived from tree biomass
through biomass models and turnover rates, parts of the input such as the litter of
the understorey are missing from tree-based products by construction, and the input
carries year-to-year variation that periodic inventories cannot capture. Uncertainty
analyses of the Finnish soil carbon estimate have identified the litter input, together
with the initial state of the soil, as the dominant sources of uncertainty
\citep{Peltoniemi2006,LehtonenHeikkinen2015}. Those analyses propagated the
uncertainty forward, by sampling the input from assumed error distributions and running
the soil model with each sample. We instead estimate the error of the input from the
soil observations themselves, together with the parameters of the models, so that the
correction applied to the input and the uncertainty of that correction are both
learned from the data and carried into every result.

The correction is a single multiplier, $\sigma_{\mathrm{input}}$, applied to the litter
input of the product over the whole simulation. The prior of $\sigma_{\mathrm{input}}$
is log-normal with a centre of $1.08$, the ratio of the total litter input to Finnish
upland soils estimated by \citet{LehtonenHeikkinen2015}, which includes the
understorey, to the tree litter input of the product used here, $2.70$ against $2.511$
tonnes of carbon per hectare per year, and a width of $0.25$ on the logarithmic scale.
Because the multiplier is one number for the whole country, it corrects the national
level of the input and not the input of individual plots, whose errors are carried by
the error model described below. To test how far the result depends on the centre of
this prior, the calibration was repeated with the centre at $1.27$, the ratio implied
by the reconstruction of \citet{Liski2006}.

The same reasoning is extended to the starting state. The litter input of 1917, which
sets the state of the soil at the start of the simulation, is not known either, and is
estimated through a second parameter, $\sigma_{\mathrm{init}}$, the ratio of the litter
input of 1917 to the litter input of 1985. The prior of $\sigma_{\mathrm{init}}$ is
log-normal with a centre of $0.90$, the ratio of the growing stock of Finnish forests
in 1917 to that in 1985 \citep{Korhonen2024} converted to a ratio of litter input with
the relation between litter input and growing stock observed in the product, and the
same width of $0.25$; the reconstruction of \citet{Liski2006}, based on the same
growing-stock record, implies a similar ratio, about $0.86$. Both parameters are estimated together
with the model parameters, so that the uncertainty of the input and of the starting
state propagates into the initial state, into the simulated trajectory and into the
predictive intervals of every model.

\subsection{Transient initialisation}

Rather than starting each plot from a soil in balance with its present litter input,
each draw of the parameter vector starts the soil in 1917 and runs it forward to the
first soil observation before the comparison with data begins. The year 1917 is a
convention, chosen as the year of Finnish independence and therefore the
point from which the national forest record is usually counted (any year starting far enough before the observations for the
soil to carry a history would work). 
The state of the soil in 1917 is obtained conventionally, with a steady state assumption with a constant litter input equal to the litter input of
1917, and the litter input is then raised along a prescribed shape over the sixty-eight
years from 1917 to 1985, after which the simulation is driven by the annual litter and
climate of the record.

The shape along which the litter input is raised is taken from the reconstruction of the litter input to Finland's
forest soils published by \citet{Liski2006}, who estimated that input annually from
1922 onwards by combining national forest inventory data with biomass and turnover
models. 
% ROUND 5, 2026-09-23 (notes 1235, 1241, 1254): area stated descriptively; both
% endpoints named (1985 end = product mean over 1985-1989 x sigma_input, as in
% J_t0_mean, N_PREINIT_SMOOTH = 5); the flux-vs-volume paragraph cut and its bias
% moved to Limitations. Natural mortality: kept in the shape; removing it moves the
% mean shape 0.450 -> 0.448 and no year by more than 0.037 (liski_preinit_shape()).
Three decisions were made about that reconstruction.
First, only the tree-derived terms are used, meaning tree litter together with harvest
residues and the litter of trees that died naturally, because those are the terms
closest to the litter product used after 1985; that product represents natural
mortality only when a death appears as a loss of basal area, but natural mortality is
between $2$ and $4$ per cent of the tree-derived input in the reconstruction, and
leaving it out changes the shape used here by less than $0.04$ in any year. The ground vegetation is left out of the shape for the
same reason that it is left out after 1985, assuming that it is absorbed by the input
multiplier described in Section~\ref{sec:inputunc}, and including it before 1985 but not after would make that
multiplier mean two different things on either side of the join. 
Second, the national totals are converted to a per-hectare basis by dividing by the
area of upland forest soil. \citet{Liski2006} report the carbon of upland soils both
as a national total and per unit area, and the ratio of the two gives an area of
$13.90$ million hectares in 1922 and $15.22$ million hectares in 2004. We placed the
growth between the two values in 1965--1980, the period in which \citet{Korhonen2024}
locate most of the expansion of productive forest land.
Third, only the shape is taken and not the level: the series is rescaled to run from
zero in 1917 to one in 1985, so that the reconstruction supplies only the path between
the two endpoints. The two endpoints are set by the two parameters described in Section~\ref{sec:inputunc}: the input at the end of the
path is the litter input of the product, averaged over 1985--1989 for each plot and
scaled by $\sigma_{\mathrm{input}}$, and the input at the start of the path is a
fraction $\sigma_{\mathrm{init}}$ of the input at the end.


The construction of the starting state is shown in Figure~\ref{fig:init}.

\hfig{F4_initialization}{0.90}{init}{\textbf{How the starting state of the soil is obtained.}
Rather than starting each plot from a soil already in balance with its present litter input, the
simulation starts in 1917 from a soil in balance with the litter input of 1917, and the litter input
is then raised over the following sixty-eight years along a shape taken from the reconstruction of
Finland's forest litter input by \citet{Liski2006}, after which the simulation is driven by the
annual litter and climate of the record. The litter input of 1917 is not known and is estimated as
part of the calibration.}


\subsection{Keeping the models as comparable as possible}\label{sec:comparable}

% ROUND 5, 2026-09-23 (notes 1362, 1367, 1400, 1407, 1410, 1413): recast from "equal
% external information" (a condition we cannot meet) to a difficulty reduced by one
% compromise per parameter class. Facts from Prior_specs/*_priors.R: simple models
% borrow Yasso07's climate centres AND widths; fractions 0/1/2 in SP1/TP2/TP3 (not 12),
% logit SD 0.4 everywhere, simple-model centres on ICBM h; sigma_init/sigma_input same
% centre (0.90/1.08) and width (0.25) in all six. The promise to report the rate
% posterior width "as a diagnostic" was dropped (never reported). Sentence 1400 dropped.
The six models cannot be given the same external information. The models differ in
structure and in the number of their parameters, and the published literature provides
different information for each of them: the three versions of Yasso come with their
own calibrations, whereas the three simple models have no calibration for forest soils.
The more the models differ in the information they are given, in their assumptions and
in their freedom, the less a difference between their results can be attributed to
their structure. We therefore reduced that difference as far as the models allow,
through one compromise for each class of parameter, summarised in the table at the end
of this section. The compromises reduce the asymmetry between the models without
removing it, and what remains is stated with each compromise.

\subsubsection{Decomposition rates}

The rates of the three versions of Yasso are fixed at their published values, for the
reasons given in Section~\ref{sec:models}. The three simple models have no rates of
their own for forest soils, and take theirs from the ICBM model
\citep{AndrenKatterer1997}: the fast rate is fixed, and the slow rate or rates are
given a very informative prior centred on the ICBM value, so that the data can move
them only slightly. The slow rates are pinned rather than fixed because the ICBM values
cannot be transferred to a forest soil unchanged. ICBM was calibrated on an arable
long-term experiment, and cropland and forest soils differ in what controls
decomposition: cultivation breaks up the soil aggregates that protect organic matter
physically in undisturbed soils \citep{Six2002}. Moreover, a rate is only meaningful
together with the climate response it was estimated with, and the simple models combine
the rates of ICBM with the climate response of Yasso07; the rates were rescaled once so
that the two climate responses agree at the climate of the experiment on which ICBM was
calibrated, but the two responses differ elsewhere. Fitting the rates freely, on the
other hand, would give the simple models a freedom that the versions of Yasso do not
have, and a narrow prior around the ICBM value is the compromise between the two.

\subsubsection{Transfer fractions}

The fractions that describe how carbon passes between pools are free in every model
that has them: twelve in each version of Yasso, and the fraction of decomposed carbon
passed on to the slower pool in the simple models, of which SP1 has none, TP2 one and
TP3 two. All fractions carry a prior of the same width, $0.4$ on the logit scale,
centred on the published value of each model, which for the simple models is the
humification coefficient of ICBM.

\subsubsection{Climate response and size of woody litter}

The three versions of Yasso take the centres and widths of their climate priors from
their own published calibrations \citep{Tuomi2009,Tuomi2011woody}, and \gap{Yasso15/20
source: Yasso20 = Viskari et al.\ 2022; Yasso15's calibration is unpublished (Järvenpää et
al., in preparation) -- decide whether to cite the FMI parameter files or the manuscript}.
The simple models, which have no published climate response of their own, borrow the
climate response of Yasso07, centre and width. Only the versions of Yasso describe how
decomposition depends on the size of woody litter, with priors taken from their own
calibrations; the simple models have no such term.

\subsubsection{The two auxiliary parameters}

Every model carries the two parameters described in Section~\ref{sec:inputunc},
$\sigma_{\mathrm{input}}$ and $\sigma_{\mathrm{init}}$, with the same prior, the same
centre and the same width, in every model.

\gap{Parameter-class by prior-criteria table (tab:priors): for each class, the centre
source, the width source, the transform and constraint, and which models carry it. It
already exists in the methods documentation.}

\subsection{Numerical integration}

% ROUND 5, 2026-09-23 (notes 1354-1377): rewritten for a reader without numerical
% analysis. The section is about the FORWARD run, not the steady state. Yasso's
% matrixexp is FMI's Yasso15 core (Jarvenpaa) with Taylor terms 10 -> 20 and a cap on
% scaling steps (ours); only the method is cited. Annual-step check:
% doublechecks/tp3_subannual_forcing_test.R.
Each model is run forward one year at a time, from 1917 to the last observation, with
the litter input and the climate of each year held constant within that year. Within a
year, each pool receives its share of the litter input and of the carbon passed on from
other pools, and loses a fixed fraction of its own carbon per unit of time. The carbon
in the pools at the end of the year can therefore be obtained in two ways. The simple
way is to compute how fast each pool is changing at the start of the year and to
assume that the pool keeps changing at that speed for the whole year. The exact way is
to use the solution of the equations over the year, which for a single pool is known
in closed form: the pool approaches the value at which gains and losses balance along
an exponential curve, and never passes that value. No numerical solver of differential
equations is used: within each year the equations of all six models have a known exact
solution, and that solution is evaluated once per year. The only approximation is the
annual time step itself, and re-running the calibrated three-pool model with twelve
monthly steps a year changed the national mean stock by at most $0.7$ tonnes of carbon
per hectare in any year, about one per cent of the stock.
% 2026-09-23: doublechecks/tp3_subannual_forcing_test.R on run 20260831_162456, posterior
% median, all 456 plots: max |monthly - annual| of the national mean 0.69 (mean 71.6).

The two ways differ when a pool turns over faster than once a year. The simple way then
carries the pool past its balance value in one year and back in the next, and produces
an oscillation that belongs to the method of calculation and not to the model. The
fastest pool of the three-pool model turns over more than once a year in warm years,
so the difference matters here.

All six models use the exact solution. For several pools that pass carbon to one
another the exact solution involves the exponential of the matrix of rates and
transfers, and the models differ only in how that exponential is obtained. For the
one-pool and two-pool models the solution is written out as a formula. For the
three-pool model carbon moves in one direction only, from the fast pool to the slow
pool and from the slow pool to the humus, so the solution can also be written out as a
formula, as a combination of one exponential decay per pool. In the three versions of
Yasso carbon moves between pools in every direction, and the exponential is computed
numerically in the Fortran code of the models by the standard method of scaling and
squaring \citep{MolerVanLoan2003}.

\subsection{Mean transit time as a common measure of a model}\label{sec:mtt}

The six models differ in how many carbon pools they contain and in how carbon moves
between those pools, so their parameters cannot be compared with one another
directly. A single number can be compared, however, and that number is the mean time
a unit of carbon spends in the soil between entering it as litter and leaving it as
respired carbon dioxide. \citet{Sierra2017} call that quantity the \emph{mean transit
time}, and we use their term throughout. Much of the soil carbon literature, including
earlier work of our own, calls the same quantity the mean residence time; Sierra and
coauthors show that the term residence time has been used for at least three different
quantities and discourage its use, so we adopted their convention. At steady state the mean transit time is also equal to the total
carbon stock divided by the total input flux, which is what many authors call the
turnover time.

% ROUND 5, 2026-09-23 (notes 1409, 1447): equation dropped; the balance stock is
% computed directly (formula for SP1/TP2/TP3, one linear solve -A C = b in Yasso,
% yasso15.f90 model_steady_state), not by a spin-up; the old "without inverting the
% matrix" was wrong for Yasso. Linearity, not the unit input, is what makes MTT
% independent of the input amount. Closing "We note that..." paragraph deleted.
For a system in which each pool loses carbon in proportion to the amount of carbon it
holds, and in which the rates do not change with time, the mean transit time can be
computed from the rates and transfers of the model alone
\citep{Rasmussen2016,Sierra2017}: at steady state the mean transit time equals the
total carbon the system holds divided by the flux of carbon entering it. We therefore
give each model a litter input of exactly one tonne of carbon per hectare per year and
compute the carbon stock at which the gains and losses of every pool balance. The
balance stock is computed directly and not by running the model until it settles: for
the three simple models the balance stock is written out as a formula, and for the
three versions of Yasso the balance equations of all pools are solved together as one
linear system. The mean transit time is then the balance stock divided by the input.
Because the models are linear, the result does not depend on the amount of input
chosen: doubling the input doubles the stock and leaves the ratio unchanged. An input
of one tonne only makes the stock and the mean transit time the same number; the
quantity reported is the ratio, in years.

Two properties make the quantity useful as a basis for comparing models. First, the
mean transit time computed in this way is a property of each model alone.
Neither the litter input multiplier, nor the estimated starting state, nor any soil
carbon observation enters the calculation, so two models can be compared even though
they were fitted to the same data with different numbers of parameters, and our
calibrated parameters can be compared with the published parameters of the same model
on exactly the same footing. For the same reason, the mean transit time is not the
ratio of simulated stock to simulated input at any date of the simulation, which
reflects the imbalance of the soil as well as the model. Second, the calculation collapses a parameter vector of
between six and twenty-six numbers into one number with a physical meaning and a unit
of years, which allows models that are otherwise very different to be compared on the
same axis.

Because a decomposition model responds to climate and to litter chemistry, the mean
transit time is not a constant of the model but a property of the model evaluated
under a specific condition for these variables. We therefore fix one reference condition and use the same
reference condition. The reference climate is the mean climate of the study plots, with a mean
annual temperature of $3.4$ degrees Celsius, a temperature amplitude of $13.8$ degrees
and an annual precipitation of $591$ millimetres. The reference litter is the mean
litter composition of the study plots, normalised to a total of one, which divides the
input into $0.766$ non-woody litter, $0.219$ fine woody litter and $0.014$ coarse
woody litter, each with its own chemical composition, so that the part of each model
describing the effect of the size of woody litter is exercised as it would be in a real
simulation. The mean transit times reported here should
therefore be read as the mean transit time of each model under such conditions, not as a universal constant of the model, nor as the transit time of the real soil under changing conditions. The comparison between our calibration and the
published parameterisation does not depend on the choice of reference, because both
sides of the comparison are evaluated at the same reference.

\subsection{The error model}

% ROUND 5, 2026-09-23 (notes 1454, 1457, 1461): 0.800 kept as a declared design value
% (option C) and defended WITHOUT the calibrated residuals (circular). Numbers from the
% input bundle of run 20260831: sd of log observed stock around campaign means 0.394
% (1205 obs); sd of log plot-mean litter across the 456 plots 0.413;
% sqrt(0.394^2 + 0.413^2) = 0.571. Kurtosis ~7 from 2026-08-11.
Observations are compared with predictions on the logarithmic scale, so that the error
is proportional to the size of the stock. The total spread of the error is fixed at
$0.800$ on the logarithmic scale, about twice the spread of the observed stocks around
the national mean of each campaign, $0.39$. That spread is the error of an idealised
model that reproduces exactly the change of the national mean from one campaign to the
next and knows nothing about individual plots. Such a model would follow the
measurements together with their errors, because the measurements are all the
information it has. A process-based model adds information that the measurements do
not contain, in the litter input, the climate and the representation of decomposition,
and where the measurements are in error a process-based model is expected to depart
from them. The error term has to leave room for that departure.

Four further reasons place the realistic error above that of the idealised model. The
litter input carries errors of its own at each plot: the variation of the litter input
between plots, $0.41$, is followed only weakly by the observed stocks, and counted as
an independent error it raises the total to $0.57$. The observed errors have heavier
tails than a normal distribution. The observed sample excludes outlying plots. Part of
the measurement error, in the stone volume, in the bulk density and in the modelled
deep tail of the profile, is shared between plots and does not appear in a spread
between plots. A value at the upper end of the plausible range was preferred, because
a larger error weakens the pull of the observations relative to the priors, so that
results attributed to the observations are stated conservatively.

% ROUND 5, 2026-09-23 (notes 1471-1530): rewritten in plain words. Averaging over the
% offsets = marginalisation; the variance rule replaces "quadrature"; the factorisation
% sentence and the correlated-vs-independent comparison removed. n_eff recomputed with
% the final spreads (all four parts, 8 equal-count bands, run 20260831): 221 (without
% the campaign part 261; the August "260" omitted it). Campaign spreads prescribed
% (code: ASSUMED); first-campaign reasons verified (LM layer; 1986-1995; Kramarenko
% 2012 sect. 2.4: C below 5 cm predicted from loss on ignition on most plots).
% sigma_e uses the later-campaign spread, so the first campaign's total is 0.802; the
% distinction is deliberately not stated in the text (Lorenzo: drop it).
Observations of soil carbon are not independent of one another. Observations of the
same plot share whatever the models get systematically wrong about that plot,
observations from the same part of the country share whatever the models get wrong
about that part of the country, and observations from the same campaign share whatever
the protocol of that campaign did differently. The error of each observation is
therefore divided into four parts: a part shared by all observations in the same band
of latitude, a part shared by all observations of the same plot, a part shared by all
observations of the same campaign, and a part that belongs to the observation alone,
\[
  \log y_{ij} \;=\; \log f_{ij}(\theta) \;+\; u^{R}_{r(i)} \;+\; u^{P}_{i} \;+\; u^{C}_{j} \;+\; e_{ij},
\]
where $y_{ij}$ is the stock observed on plot $i$ in campaign $j$, $f_{ij}(\theta)$ the
stock predicted by the model, and $u^{R}$, $u^{P}$, $u^{C}$ and $e$ the four parts.
Each part is treated as a random error with a mean of zero and a fixed standard
deviation, and no part is estimated for any individual plot, band or campaign.
Instead, the likelihood is averaged over all the values each shared part could take,
weighted by how probable each value is, which for normal errors can be done exactly;
the shared parts then disappear from the calculation and leave only a correlation
between observations that share a plot, a band or a campaign.

The standard deviations of the shared parts were taken from the observations
themselves, after removing the mean of each campaign. The standard deviation of the
regional part, $0.117$, is the spread between the means of eight bands of latitude,
whose boundaries were placed so that each band contains the same number of plots,
because the plots are unevenly distributed from south to north and bands of equal width
would leave some bands almost empty. The standard deviation of the plot part, $0.396$,
is measured from how strongly two observations of the same plot vary together.
Independent errors add by their variances, that is by the squares of their standard
deviations, and not by the standard deviations themselves, so the variances of the four
parts add up to the variance of the total, $0.800^2$, and the standard deviation of the
part that belongs to each observation alone, $0.685$, is what remains once the three
shared parts are accounted for. Dividing the error in this way does not change how far
a single observation may lie from the model, but it changes how much the observations
are worth together: observations that share part of their error partly repeat one
another, and the $1205$ observations carry about as much information on the national
mean as $220$ independent observations would. Both published calibrations of Yasso
assumed the errors of the observations to be independent \citep{Tuomi2009,Viskari2022},
and the calibration of Yasso20 named the correlation of errors within a time series as
an effect still to be included \citep{Viskari2022}. Treating correlated observations as
independent leads to parameter uncertainties that are too narrow
\citep{Oberpriller2021}, and the shared parts of the error used here are a way of
avoiding that.

The standard deviations of the campaign parts cannot be measured from the data,
because a shared offset between campaigns and a genuine change of the stock between
campaigns look the same in the observations. The standard deviations were therefore
prescribed: $0.030$ for each of the two later campaigns, which followed similar
protocols, and twice that, $0.060$, for the first campaign, which differs from the
later two in three documented ways. The organic layer of the first campaign excluded
the litter and moss layer, which is added here from each plot's measurement of 2006;
the fieldwork of the first campaign extended over a decade; and the carbon content of
its mineral layers below $5$ cm was mostly predicted from the loss of mass on ignition
rather than measured directly \citep{Kramarenko}.

\subsection{Sampling}

Each model was calibrated with five independent runs of $50\,000$ iterations, of which
the first $5\,000$ were discarded, using the differential evolution Markov chain sampler
of \citet{terBraakVrugt2008}, in the implementation provided by the
\texttt{BayesianTools} package \citep{BayesianTools}, version $0.1.8$, under R $4.3.1$
\citep{RCore}.
% ROUND 5, 2026-09-23 (note 1537): "snooker update" explained in words, name dropped.
The sampler proposes each new parameter vector from the difference between past
states of its chains, occasionally along the line joining the current state to a past
one, so that the proposals adapt to the shape of the posterior; this suits posteriors
with strongly correlated parameters, such as the one described in this paper. Each run
maintains three interacting chains, so that the five runs together provide fifteen
chains.
% ROUND 5, 2026-09-23 (note 1713): KL and the nat defined, with two Gaussian anchors
% (quarter width: -ln(0.25) + 0.25^2/2 - 1/2 = 0.92; one-SD shift: 1/2). Computed as
% in calibration_engine.R kl_from_kde(): KDE of marginal samples, natural log.
The information that the observations add to each parameter was measured as the
Kullback--Leibler divergence of the marginal posterior from the marginal prior,
estimated from kernel densities of the samples. The divergence is zero when the
posterior is identical to the prior, that is when the observations say nothing about
the parameter, and grows as the observations narrow the distribution or move the
distribution away from the prior. The divergence is expressed in nats, the unit that
results when information is measured with the natural logarithm, as bits are the unit
that results with the logarithm to base two. Two cases give a sense of the scale: a
posterior with the same centre as its prior and one quarter of its width corresponds to
about $0.9$ nats, and a posterior with the same width as its prior but moved by one
prior standard deviation corresponds to $0.5$ nats. Convergence
was assessed with the potential scale reduction factor and with the effective
sample size, and the rate of proposals rejected as numerically impossible was
monitored throughout.

\subsection{Predictive evaluation and residual analysis}

Predictions were generated for both the calibration plots and the held-out plots by
drawing from the posterior and running each draw forward. Model skill is reported
as the coefficient of determination and the root mean squared error, separately for
the calibration plots and for the held-out plots. Residuals, defined as the
logarithm of the observed stock minus the logarithm of the predicted stock, were
then regressed against site variables, and the importance of each site variable was
assessed with a random forest evaluated out of bag.

\subsection{Forward simulations}\label{sec:forward}

% ROUND 5, 2026-09-23 (Conclusions notes 2139, 2149): side analyses, "an order of
% magnitude, not a forecast". Code: Calibration_real_data_transient/forward_scenarios.R
% (inside the six predictive scripts, same draws; existing outputs verified
% byte-identical) and manuscript/figures/build_forward_scenarios.R. Rates over WHOLE
% 20-year climate cycles: decadal rates alternate with the replayed 2005-14 / 2015-24
% climate (doublechecks/yasso20_late_source_check.R).
Two side analyses extend the calibrated models beyond the observations, to give an
order of magnitude for the accumulation that remains rather than a forecast. Both use
the same $100$ posterior draws as the predictions and the plots measured in all three
campaigns. The projection runs each model for $60$ years after the last observation,
with the litter input held at its value in the last observed year, 2024, and with the
climate of the last $20$ observed years repeated three times, so that the projection
contains neither a change of climate nor a change of forest growth. Because the
repeated climate alternates between a cooler and a warmer decade, rates of change over
the projection are computed over whole $20$-year cycles. The equilibrium stock of each
plot is the stock at which gains and losses would balance if the same litter input and
climate were held indefinitely, computed with the routine of each model for the state of
balance, as for the mean transit time (Section~\ref{sec:mtt}); the distance from balance
is the equilibrium stock divided by the stock of 2024, minus one. The input scenario
repeats the projection with the litter input increased by $20$ per cent from the first
projected year onwards. Because the models are linear, the eventual gain of stock under
that scenario is $20$ per cent of the equilibrium stock, and the scenario shows how much
of the gain is realised within $60$ years.

% =============================================================================
% RESULTS AND DISCUSSION. Combined, and nested at the subsection level: the first
% paragraph of each subsection states what was measured and stops; the paragraph
% that follows interprets it and opens by naming what it is interpreting.
% =============================================================================
\section{Results and discussion}

% ROUND 5, 2026-09-23 (notes 1575-1732): the five opening subsections ("reach the
% observed stocks", "bracket the observed rate", "shape of the trajectory", "structure
% nearly invisible", "the data fix a stock") collapsed into one short subsection: the
% fit is a precondition, not a result. Old text in git (commit dcd1b88). The mechanism
% of the transient start (notes 1664, 1669) is deliberately NOT claimed here: the
% balanced-start counterfactual has not been run. F5 moved to the supplement.
\subsection{The models reproduce the national trajectory, and the fit cannot tell them
apart}\label{sec:fit}

All six calibrated models reproduce the mean soil carbon stock of the three campaigns
to within $3.5$ tonnes of carbon per hectare (Figure~\ref{fig:modelmean}), and over
the full observed window the six models bracket the observed rate of change of
$+0.259$ tonnes of carbon per hectare per year, from below in SP1 and Yasso20 and from
above in TP2 and TP3. Agreement of this kind is expected, because the three campaigns
are the data the models were calibrated on, and it shows consistency with the
observations rather than predictive ability \citep{LeNoe2023}. The observations also
carry their own uncertainty: the three campaigns followed three protocols and the
first campaign spans a decade, so that the rate over the full window survives a shared
offset between campaign levels of about $9$ per cent, whereas the rate over the last
interval, on which the apparent slowing of the accumulation rests, is removed by an
offset of about $2$ per cent. Which model lies closest to the observed rate cannot
therefore be read as evidence that the model is better.

\hfig{F3_mean_soc}{0.92}{modelmean}{\textbf{Modelled and measured mean soil carbon, 1985--2024.} Lines are the mean stock across plots for each of the six models, and the shaded bands are $95\%$ posterior intervals. Markers are the measured campaign means with their confidence intervals, calculated on the same set of plots as the model curves. All six models pass within $3.5$ tonnes of carbon per hectare of every campaign mean, and the shape of the modelled trajectories separates the six models into three groups.}

\suppfig{F5\_stock\_change}

The fit does not separate the model structures either. Skill does not increase with
the number of pools: the coefficient of determination is $0.016$ to $0.021$ on the
calibration plots and no more than $0.0012$ on the held-out plots, and the root mean
squared error on the held-out plots differs by $3.6$ per cent between the best and the
worst model (Table~\ref{tab:models}). Skill is measured plot by plot, so the
coefficient of determination measures how well the models separate one plot from
another, and none of the six represents the local variation that separates them
(Section~\ref{sec:plotscale}); the national trajectory, which all six reproduce, is a
different quantity. Inside the models, the observations leave most parameters
undetermined: all parameters converged, but the transfer fractions stay close to their
priors, gain less than one nat of information, and trade off against the other
fractions leaving the same pool. The observations determine the national level and
trend of the soil carbon stock, and not how each model produces them.

\begin{table}[htbp]\centering\small
% auto-generated by build_T1_model_summary.R -- per-model summary
\begin{tabular}{lccccccc}
\hline
Model & Pools & Free & $R^2$ & $R^2$ & RMSE & Eff.\ flux & $\sigma_{\mathrm{init}}$ \\
 & & params & (cal) & (hold) & (hold) & $\sigma_{\mathrm{inp}}\bar J$ & \\
\hline
SP1 & 1 & 6 & 0.017 & 0.000 & 46 & 4.9 & 0.67 \\
TP2 & 2 & 7 & 0.018 & 0.000 & 47 & 6.4 & 0.40 \\
TP3 & 3 & 9 & 0.019 & 0.000 & 47 & 6.5 & 0.37 \\
Yasso07 & 5 & 20 & 0.028 & 0.002 & 45 & 6.6 & 0.24 \\
Yasso15 & 5 & 26 & 0.031 & 0.004 & 44 & 5.9 & 0.18 \\
Yasso20 & 5 & 26 & 0.023 & 0.001 & 45 & 6.6 & 0.21 \\
\hline
\end{tabular}

\caption{\textbf{Summary of the six calibrated models.} Number of carbon pools, number of free
parameters, coefficient of determination on calibration plots and on held-out plots, root mean squared
error, the effective litter flux implied by the calibration, and the estimated ratio of the litter
input of 1917 to the litter input of 1985. Skill does not increase as pools are added.}
\label{tab:models}
\end{table}

\suppfig{F6b\_skill\_vs\_complexity}
\suppfig{F8b\_fraction\_ridges}

The six trajectories nevertheless differ in shape, and the shape is what an
extrapolation rests on. Yasso07 and Yasso15 rise and then flatten, TP2 and TP3 keep
rising, and SP1 and Yasso20 reach a maximum in the late 2000s and then decline. Models
that agree within the observed window can therefore disagree outside it.

\subsection{The calibration settles at a larger litter input, and in two versions of
Yasso at a faster turnover, than the published parameterisations}
% ROUND 5, 2026-09-23 (P1, note 2114): "requires" removed throughout. Fit cost of the
% published MTT from F14_mrt_ridge.rds (run 20260831, built 2026-09-11): best ll within
% +-5% of the published POINT vs best overall = 9.7 / 3.0 / 0.0 (upper bounds);
% P(MTT >= published) = 0.2% / 6.5% / 61%.

% ROUND 5, 2026-09-23 (note 1732): no section split; the collapsed opening subsection ends on the hinge sentence.

The calibrated models sit away from their published parameterisations in two
quantities at once, and the displacement in transit time is graded across the three
versions of Yasso rather than uniform. The effective litter input, defined as the
multiplier $\sigma_{\mathrm{input}}$ times the litter input carried by the product,
ranges from $3.19$ to $4.34$ tonnes of carbon per hectare per year, against the
$2.511$ tonnes that the product supplies. The intrinsic mean transit time is $24.6$
years for Yasso07, $27.1$ years for Yasso15 and $22.6$ years for Yasso20, against
published values of $33.5$, $30.4$ and $21.9$ years respectively. Expressed as a
ratio of the calibrated transit time to the published transit time, the three
versions of Yasso give $0.73$, $0.89$ and $1.03$.

Yasso20 therefore serves as a negative control, and the negative control is worth
stating explicitly. The calibration displaces the transit time of Yasso07 strongly,
displaces the transit time of Yasso15 moderately, and does not displace the
transit time of Yasso20 at all: the calibrated value for Yasso20 lies three per cent
above the published value, and the $90$ per cent posterior interval of the calibrated
value, from $19.6$ to $26.4$ years, contains the published value comfortably. A
calibration that moved every model in the same direction by a similar amount would be
open to the objection that the movement came from the data treatment or from the
likelihood rather than from the models. A calibration that moves two models and leaves
the third where its authors put it is not open to that objection.

All transit times quoted here are compared against the published \emph{point}
parameterisation. Yasso15 and Yasso20 also publish posterior samples, and because
transit time is a non-linear function of many parameters, the median transit time
of a published sample is not the transit time of the median published parameters.
The two comparators are never mixed in this paper. 

\hfig{F9_auxiliary_sigmas}{0.98}{sigmas}{\textbf{The two estimated auxiliary quantities, read as physical quantities.} \emph{Left:} the litter input multiplier expressed as an effective flux of litter carbon to the soil, that is, the multiplier times the litter input that the product supplies. The dotted line marks the flux the product itself supplies, and the shaded bands mark the envelope of independent estimates of boreal net primary production. Every model settles above the tree litter the product supplies and inside the envelope. \emph{Right:} the estimated ratio of the 1917 litter input to the 1985 litter input. All six models place the 1917 input below the 1985 input, which is the direction the forest history requires.}



\hfig{F12_mrt_yasso}{0.88}{mrt}{\textbf{Mean transit time of carbon in the three versions of Yasso, calibrated and published.} Transit times are computed at a common reference climate and litter composition, from the model equations alone, so that the comparison depends neither on the litter input multiplier nor on the starting state.}


The displacement in litter input and the displacement in transit time are one
displacement and not two. The observed stock constrains the product of the litter
input and the transit time, so a soil that receives more litter and decomposes it
faster fits the observations as well as a soil that receives less litter and
decomposes it more slowly. The calibration moves along that trade-off.
Moving along the trade-off back towards the published transit time costs little in
fit. Among the posterior draws that lie within five per cent of the published transit
time, the best fit is at most $9.7$ log-likelihood units below the best fit overall for
Yasso07, at most $3.0$ units for Yasso15, and no lower for Yasso20, and $6.5$ per cent
of the posterior of Yasso15 lies at or beyond the published value. The soil
measurements therefore suggest a direction rather than impose a value.


\hfig{F14_mrt_ridge}{0.98}{ridge}{\textbf{The trade-off between transit time and litter input.} For each version of Yasso, the best fit to the data attainable at each combination of mean transit time and litter input multiplier. \textbf{The colours show the fit to the data alone, with the priors excluded}, from blue at the best fit through green and yellow to red at the worst. \textbf{The white contour lines show the posterior, which combines the fit to the data with the priors.} The dashed line marks the published transit time and the white dot marks the posterior median. Cells that the sampler visited fewer than five times are left blank, and a blank cell means that the region was not visited rather than that the region cannot be fitted. The offset between the bluest colour and the innermost contour measures how far the priors pull the answer away from where the data alone would place it.}
 \suppfig{S13\_mrt\_ridge\_benchmark}

The trade-off is not created by the choice of starting state, and a transient
starting state does not remove it. The trade-off follows from the condition of
balance itself, so the trade-off exists under either choice of starting state. What
a transient starting state adds is information about the rate at which the soil
approaches its eventual stock, and the rate of approach depends on the transit
time alone. Estimating the starting state therefore narrows the trade-off, and the
narrowing is partial: across the three versions of Yasso the transit time and the
litter multiplier still correlate between $-0.64$ and $-0.73$.

\subsection{The displacement in litter input follows the data more than the prior}

% ROUND 5, 2026-09-23: "two production runs" reframed as a designed sensitivity
% analysis (no reference to development history); reproducible with
% HIKET_SIGMA_INPUT_CENTRE=1.27 (RELEASE_NOTES.md). Arm A = 20260820_1554*.
The litter input multiplier could in principle be an artefact of the prior placed on
it. To test that possibility, the calibration was repeated with the centre of that
prior moved from $1.08$, the value derived from \citet{LehtonenHeikkinen2015}, to
$1.27$, the value implied by the reconstruction of \citet{Liski2006}, with everything
else unchanged. Moving the prior centre by $18$ per cent moves the posterior median only
about one third as far: from $1.54$ to $1.62$ in Yasso15, from $1.73$ to $1.87$ in
TP2, and from $1.27$ to $1.35$ in SP1. Under the higher prior centre the posterior
median still sits between $1.28$ and $1.50$ times above the prior centre, and under
the lower prior centre the posterior median sits between $1.17$ and $1.60$ times
above it.

Because the posterior follows the prior centre only partially, and because the
posterior remains well above the prior centre under either choice, the preference for
a larger litter input comes mainly from the soil measurements rather than from the
prior. The preference survives either anchor, and choosing between
the two anchors is therefore a question about how the prior should be documented and
not a question about whether the result holds.

% sigma_init posterior discussed in Section sec:headroom (ROUND 5, 2026-09-23).

\subsection{What might be missing from the litter input}

% CARRIED in part from the previous manuscript ("the multiplier has a physical
% referent"), and extended with the literature review completed 2026-09-04.

The multiplier $\sigma_{\mathrm{input}}$ is worth reading as a physical quantity and
not as a nuisance parameter, because the multiplier is not scale-free. Multiplied by
the litter input the product supplies, the multiplier becomes an effective flux of
litter carbon to the soil, and an effective flux can be compared with independent
budgets. The gap to be explained is between $0.68$ and $1.83$ tonnes of carbon per
hectare per year.

The most likely explanation is not that the litter model is wrong but that the
litter model is incomplete. Litter products of this kind are trusted because they
have worked for decades, and one reason they have worked is that a calibration free
to adjust its decomposition rates can absorb a missing flux without any visible
misfit. A calibration performed on externally anchored decomposition rates removes
that compensation, and a missing flux then has to appear somewhere.

\subsubsection{Carbon released by living roots}

Carbon released into the soil by living roots belongs to no biomass compartment, so
a litter model assembled compartment by compartment cannot represent it. Agricultural
carbon accounting has treated the same flux as a standard term for two decades. The
widely used approach of \citet{Bolinder2007} adds an explicit pool of extra-root
material set at $65$ per cent of the standing root biomass, derived as the ratio of
the share of below-ground carbon released by living roots and retained in the soil
to the share that remains in the roots themselves. The extra-root pool is defined as
``exudates, as well as root hairs and fine roots sloughed off during the growing
season which, because of sampling difficulties, are not included in the `root'
fraction'' \citep{Bolinder2007}. Nothing of the kind appears in the forest inventory
chain.

Carbon released by living roots also forms stable soil organic matter more
efficiently than above-ground litter does. In the Ultuna long-term experiment the
humification coefficient fitted for root-derived carbon was about $2.3$ times the
coefficient fitted for the same mass of above-ground crop residues
\citep{Katterer2011}. A below-ground input that is both unaccounted for and
disproportionately efficient at forming stable organic matter is a candidate of the
right kind for a displacement inferred from stocks.

The agricultural coefficient cannot be transferred to forests unchanged. In an
annual crop the root biomass is measured once, at harvest, so the coefficient of
$65$ per cent compensates for everything the single measurement misses, including
the roots that grew and died during the season. A forest litter chain already
represents the turnover of fine roots explicitly, through a turnover rate applied
each year, so the forest equivalent of the extra-root term covers less. The forest
equivalent covers more than nothing, however, because forest fine-root biomass is
itself estimated from soil cores, and soil cores miss sloughed material for exactly
the reason \citet{Bolinder2007} give.

\subsubsection{Understorey vegetation}

The litter product used here represents tree litter and does not represent
understorey vegetation, by design. The national inventory does represent understorey
vegetation, and values it at $0.506$ tonnes of carbon per hectare per year in
southern Finland and $0.666$ in northern Finland
\citep[Table 6.4-3, after][]{StatFin2024NID,MuukkonenMakipaa2006}. The same values
appear in the appendix of \citet{LehtonenHeikkinen2015}. Understorey vegetation is
therefore both the best-quantified of the candidate fluxes and a flux that the input
used here is known to omit. 

\hfig{SX_input_vs_literature}{0.98}{inputlit}{\textbf{The litter input inferred by the calibration, against independent estimates of the same flux.} Green symbols are inventory-based litter products, blue symbols are independent estimates of net primary production, which bound the litter flux from above, and the lower group shows the effective flux inferred here for each of the six models. Bars are reported ranges for the literature estimates and $95\%$ posterior intervals for the models.}


The share of the understorey rises towards the north, so a single national
multiplier applied to a tree litter product must be wrong in opposite directions at
the two ends of the country. \citet{Lehtonen2016} reached the same conclusion from
the other direction, reporting that ``the role of understorey litter input was
underestimated when the Yasso07 model was parameterised, especially in northern
Finland''.

\subsubsection{Mycorrhizal fungal residues}

Where boreal soil carbon comes from has been revised substantially over the last
decade. Using bomb radiocarbon across a chronosequence of forested islands,
\citet{Clemmensen2013} attributed $50$ to $70$ per cent of the stored soil carbon to
roots and to root-associated organisms rather than to above-ground litter, and
identified mycorrhizal mycelium as a central agent of the storage. Production of
mycelium in coniferous stands is measured in hundreds of kilograms per hectare per
year \citep{Ekblad2013,Hagenbo2017}, of which only a part becomes input to the soil.

Mycelium is not a tree biomass compartment, so no turnover coefficient in a litter
chain can produce it, and the flux is missing by construction rather than
underestimated. The figure of $50$ to $70$ per cent should not be read as the size
of the missing flux, because the figure describes the origin of carbon that has been
stored and part of that carbon already enters our input as fine-root litter. What
can be said is that the dominant precursor of stable boreal soil carbon is
root-derived and partly fungal, and that the fungal part is outside inventory litter
models.

\subsubsection{Fine root lifespan, where the evidence points the other way}

Fine roots are the largest below-ground litter term, and fine-root litter enters the
inventory chain as a biomass multiplied by a single turnover rate of $0.85$ per year
applied across species and regions. The lifespan behind that rate is among the least
settled quantities in boreal ecology. Direct observation with minirhizotrons puts the
median lifespan of Norway spruce fine roots below two and a half years, whereas
radiocarbon dating of the same species returns apparent ages of up to fourteen years
\citep{Strand2008}.

A reconciled mean lifespan of about $3.8$ years would imply a turnover rate near
$0.26$ per year, which is roughly a third of the rate in use, and a lower turnover
rate would make fine-root litter smaller rather than larger. The fine-root evidence
therefore points in the opposite direction to the other three candidates, and would
widen rather than close the gap. We report the fine-root evidence as it stands. A
section in which every candidate mechanism happened to support the conclusion would
be less credible than one in which a candidate does not.

\subsubsection{Early-succession vegetation}

Litter models keyed to standing tree biomass are least informative in the years
immediately after a stand-replacing harvest, because tree biomass is then near zero
while the site is neither bare nor unproductive. Ground vegetation biomass falls
sharply at clear-cutting and then recovers through a pioneer phase in which a large
share of annual production returns to the soil each year \citep{Palviainen2005}.
Whether the total over a rotation exceeds what a stand-keyed model assumes is not
established, and we raise early-succession vegetation as a gap in the timing of the
input rather than as a quantified addition to its total.

\subsubsection{What the candidates add up to, and what they do not}

Two of the candidate fluxes can be given Finnish numbers. Understorey vegetation
supplies between $0.51$ and $0.67$ tonnes of carbon per hectare per year, and
mycorrhizal residues plausibly supply a further $0.2$ to $0.5$. The two together
reach the lower end of the gap of $0.68$ to $1.83$. Carbon released by living roots,
priced at the agricultural coefficient and applied to the fine-root share of our own
litter product, would add a further $0.73$, which taken at face value would carry the
total past the upper end of the gap.

The candidate fluxes should therefore be read as a set that brackets the gap and not
as a sum that matches it. A budget that added up exactly would be the more suspicious
result, because the candidate fluxes overlap, because the agricultural coefficient
must be smaller in forests, and because one of the candidates points the other way.
The calibration settles at a larger litter input than the product supplies, several
physically plausible sources of the difference exist, and identifying which of them is
responsible is not a task this study can complete.

\subsection{Does it matter where the calibration lands on the trade-off?}

The soil measurements cannot choose a point on the trade-off between litter input
and transit time, so the choice has to be judged by its consequences. If moving
along the trade-off changes nothing that can be observed, then the displacement
reported above is bookkeeping. If moving along the trade-off changes the
projection, then the displacement matters to anyone who uses these models to
project.

We compared the published parameterisation and our calibration by running both
forward under the same warming scenarios, without recalibrating either, and with
only the litter input multiplier refitted so that both reach the observed window
equally well. Because both arms reproduce the observed window, any divergence
after the observed window is extrapolation and not fit. 

\hfig{F15_forward_arms}{0.90}{forward}{\textbf{Projections from the published parameterisation and from the calibration, under the same scenarios.} Neither arm was recalibrated, and only the litter input multiplier was refitted, so that both arms reproduce the observed window equally well and any later divergence is extrapolation rather than fit. Only Yasso07 separates the two arms.}


Only Yasso07 separates the two arms. In Yasso07 the two arms differ by $-2.08$
tonnes of carbon per hectare, with an interval from $-3.57$ to $-0.79$, under the
SSP2-4.5 scenario. Yasso15 and Yasso20 do not separate at all.

The reason only Yasso07 separates is structural. Yasso07 applies a single climate
modifier to every carbon pool, so the turnover of the soil and the sensitivity of
the soil to temperature are governed by the same parameters, and moving along the
trade-off moves both together. Yasso15 and Yasso20 apply separate climate modifiers
to different pools, and the pools that hold most of the carbon respond almost
identically in the two arms. \suppfig{A1\_forward\_appendix}

The result about Yasso15 and Yasso20 is not that they are right. The result is that
Yasso15 and Yasso20 are insensitive to a calibration choice that the available data
cannot settle, which is a narrower statement than saying that our calibration
changes the projection, and a more useful one for an inventory.

\subsection{Holding or increasing the litter input}\label{sec:headroom}

% ROUND 5, 2026-09-23: numbers from manuscript/figures/forward_scenarios.rds (run
% 20260831, 100 draws, 310 balanced plots), medians. Headroom SP1 +0.6, TP2 26.1, TP3
% 27.4, Yasso07 27.3, Yasso15 41.6, Yasso20 34.7 %. Sink per 20-yr cycle, first/last:
% TP2 .222/.129, TP3 .209/.126, Y07 .106/.048, Y15 .098/.054, Y20 .071/.045, SP1 ~0.
% +20%: extra 2050 6.1-10.2 (TP3..SP1), 2084 8.9-13.3; share of eventual gain by 2084
% Y07 .55, Y15 .49, Y20 .46, TP2 .55, TP3 .54, SP1 .94.
With the litter input held at its present value, five of the six models place the
soil well below its state of balance: the equilibrium stock lies $26$ to $42$ per cent
above the stock of 2024 in TP2, TP3 and the three versions of Yasso ($27$ per cent in
Yasso07, $42$ in Yasso15, $35$ in Yasso20). SP1, whose single pool turns over slowly and
uniformly, is within one per cent of balance. \citet{Liski2006} estimated that Finnish
upland soils would eventually settle at a level $38$ per cent above that of 1922, a
value inside the range of the three versions of Yasso. The two estimates were obtained
by different routes, a reconstruction started in balance and a calibration started out
of balance, and agree in magnitude although their reference years differ.

The distance from balance is realised slowly. Over the first $20$ years of the
projection the sink is $0.07$ to $0.11$ tonnes of carbon per hectare per year in the
three versions of Yasso and $0.21$ to $0.22$ in TP2 and TP3; over the last $20$ years it
has fallen to about $0.05$ in the versions of Yasso and $0.13$ in TP2 and TP3, and by
2084 the versions of Yasso have realised $14$ to $21$ per cent of their distance from
balance. None of the six models turns into a source. The sink declines within decades
because the pools that turn over fastest approach balance first, whereas most of the
remaining distance lies in the slowest pool, which approaches balance over centuries
\citep{Sierra2018}. Under a constant litter input, the sink therefore saturates long
before the stock does.

% ROUND 5, 2026-09-23: sigma_init as an effective starting deficit. Within-model
% correlations over the 100 draws (sigma_init vs headroom -0.79..-0.97; vs sink 2025-44
% -0.89..-0.95); 1917 input / product 1985 = median of sigma_init*sigma_input over the
% full posterior, 0.83-0.97 in five models (SP1 1.07); rise 1/sigma_init x1.70-2.05.
The remaining sink is inherited from the estimated starting state. Within each model,
the parameter $\sigma_{\mathrm{init}}$ correlates with the distance from balance at
$-0.79$ to $-0.97$ and with the projected sink of the first $20$ years at $-0.89$ to
$-0.95$ across the posterior draws: the lower the litter input of 1917 relative to that
of 1985, the more the soil still has to gain. The calibration places
$\sigma_{\mathrm{init}}$ at $0.49$ to $0.60$ in five of the six models, with upper $90$ per
cent limits of $0.64$ to $0.76$, whereas the growing-stock record of the national forest
inventory implies about $0.90$ and the reconstruction of \citet{Liski2006} about $0.86$;
only SP1, at $0.83$, is compatible with that record. In absolute terms the calibrated
litter input of 1917 is not low, between $0.83$ and $0.97$ times the input the product
supplies for 1985 in the five models, because $\sigma_{\mathrm{input}}$ raises it; what
departs from the record is the rise of the input between 1917 and 1985, a factor of $1.7$
to $2.1$ in the calibration against about $1.1$ in the growing stock. Of the two halves
of the starting state, the ratio of inputs rests on a growing stock that the forest
inventory has measured since the 1920s, whereas the balance of the soil of 1917 with its
own input rests on no measurement at all, and the centuries of slash-and-burn
cultivation, tar burning and forest grazing described in the Introduction argue against
it. A soil depleted by that history would lie below balance with its own
input, and a single parameter can express such a deficit only as a lower input. We
therefore read $\sigma_{\mathrm{init}}$ as an effective starting deficit of the soil
rather than as a ratio of litter inputs. On that reading, the soil of 1917 was more
depleted than its litter input alone would imply, the remaining sink is conditional on
that depletion, and a soil assumed to start in balance with its present input would have
almost nothing left to gain under a constant input, as SP1 illustrates.

Increasing the litter input by $20$ per cent from 2025 raises the stock by $6$ to $10$
tonnes of carbon per hectare by 2050 and by $9$ to $13$ tonnes by 2084, relative to the
projection with a constant input, which corresponds to an additional sink of $0.2$ to
$0.4$ tonnes of carbon per hectare per year averaged over 2025--2050. Within $60$ years
the versions of Yasso realise about half of the eventual gain, $46$ to $55$ per cent. The
gain is larger than the sink of the projection with a constant input, but the gain is
finite, largest in the first decades, and lasts only as long as the larger input is
maintained. Neither projection contains a change of climate, and neither is a forecast.

\subsection{The models reproduce the national mean but not the differences between
plots}\label{sec:plotscale}
% ROUND 5, 2026-09-23 (note 1987): retitled; the opening paragraph reconciles national
% skill with the absence of plot-level skill (local ecological variation the models
% are not given) and hands over to the residual subsection.

The same models that reproduce the national mean trajectory have no skill in
distinguishing one plot from another, and the two results are not in contradiction.
The national mean depends on the total input of litter and on the rate at which the
soil turns that input over, which are the two quantities the calibration constrains.
The difference between two plots depends in addition on local ecological variation,
in the history of each plot and in the properties of its soil, and on how accurately
the litter input of each plot is known, and none of the six models is given that
information: all plots share one estimated initial state, no model has a term for the
properties of the soil, and the litter input of each plot comes from a product keyed
to its standing trees. Errors of that kind largely cancel in a mean over several
hundred plots and remain in full on each individual plot. The national inventory
reports national and regional totals, for which the first result is the relevant one,
and what the models miss between plots is not noise: as shown in
Section~\ref{sec:residuals}, about half of it can be recovered from site variables
that are already available.


Plot by plot, none of the six models has any skill. The coefficient of
determination on held-out plots is zero to three decimal places in all six models.
The informative part of the comparison is not the absence of skill but the
structure of what the models predict. The spread of the predicted stocks is wider
than the spread of the observed stocks, and the predicted stocks are ordered by the
standing basal area of the plot, whereas the observed stocks are not.


\hfig{F6_obs_vs_pred_holdout}{0.98}{plotscale}{\textbf{Predicted against measured soil carbon on held-out plots.} Points are coloured by the standing basal area of the plot, in quintiles. The dashed line is the one-to-one line, the dotted line is the prediction of the mean, and the solid line is the fitted relation. The predicted stocks spread more widely than the measured stocks and are ordered by basal area, whereas the measured stocks are not.}


Measured directly, the slope of the logarithm of the stock against basal area is
between $+0.044$ and $+0.049$ per square metre per hectare in the six models, and
$+0.0044$ in the observations. The residual therefore falls with basal area in all
six models, with slopes between $-0.032$ and $-0.044$ and with $p$ values from
$10^{-16}$ to $10^{-111}$, on the calibration plots, on the held-out plots, and on
the two pooled. Adjusting the observed relation for region and for soil type reduces
the observed slope further, to $+0.0018$, at which point the observed relation is no
longer statistically significant.

The gradient does not originate in the models. The six models contain no term in
basal area at all. What the six models contain is the litter input, and the litter
input itself rises with basal area at a slope of $+0.072$, which is a span of about
eighteen-fold across the observed range of basal area. The modelled stocks rise with
basal area at a slope of $+0.073$ to $+0.076$. The models therefore transmit the
gradient of the litter product essentially unchanged, and the forward simulation
then damps the gradient rather than amplifying it.

We tested whether the discrepancy could instead be a mismatch between compartments.
The stock simulated by Yasso covers dead organic matter as well as soil organic
matter, and woody litter enters the simulated stock, whereas the observations are
soil measurements that cannot contain a stump. Setting the woody litter inputs to
zero moves the modelled gradient by six per cent, from $0.0755$ to $0.0710$ in
Yasso15, against a fifteen-fold difference to the observed gradient. The three
simple models provide the decisive control, because the three simple models carry a
single total litter flux with no information about composition at all, and the three
simple models produce the same gradient. A mechanism that is absent from three of
the six models cannot explain a gradient that is present in all six.

Because the gradient is transmitted rather than generated, the discrepancy is a
statement about the litter product and about the soil, and not about the
decomposition models. The litter product states that the carbon input of a plot
varies about eighteen-fold across the range of basal area. The soil states that the
carbon stock of a plot varies by a factor of about $1.2$ across the same range. Two
readings follow, and the two readings are not exclusive.

The first reading is that the carbon stock of an individual plot is not set by the
current litter input of that plot. Each plot sits at a different point on its own
history of disturbance, and one initial state estimated for the country as a whole
cannot express what makes individual plots differ from one another. The first
reading is the argument of this paper applied one level down.

The second reading is that an eighteen-fold gradient in litter input between plots
is not credible. A litter product keyed to standing tree biomass gives a freshly cut
plot almost no input, and a freshly cut plot does not stop receiving carbon: ground
vegetation, dying roots and harvest residues continue to supply it. The residuals
support the second reading directly. Plots that were cut between 1985 and 1995 hold
about $40$ per cent more carbon than the models predict for them, relative to plots
that were not cut, and the residual falls with stand age. The failure at plot scale
is therefore evidence for the same missing flux as the national-scale result,
reached from the opposite direction, and it points at the two components of the
input argument that are otherwise weakest: understorey vegetation and
early-succession vegetation.

\subsection{Most of what the models leave behind can be recovered from site
variables}\label{sec:residuals}

About half of the residual variance is predictable from site variables that are
already available. A random forest fitted to the residuals explains between $44$ and
$65$ per cent of what the models leave, evaluated out of bag, against a repeatability
of the observations of about $0.49$. 

\hfig{F11_rf_heatmap}{0.86}{residuals}{\textbf{Site variables that predict what the models leave behind.} Relative importance of each site variable in a random forest fitted to the model residuals, scaled within each model so that the most important variable equals one. The same variables lead in all six models.}


Because the recoverable share is larger than the noise of the observations, what the
models leave behind is not a ceiling set by measurement error. The same site
variables lead in all six models: stand basal area first, then development class and
coarse-fragment content, with the importance spread across many variables rather
than concentrated in one. Two of the leading variables are proxies for productivity
and one is a physical property of the soil. A productivity proxy leading in every
model is what a spatially varying litter input would look like. A physical property
of the soil leading in every model is something for which none of the six models has
a term at all. We record the observation and stop there, because choosing the method
that would exploit it is not the task of this paper.

\subsection{What the result may mean for soil carbon modelling and its applications}
% ROUND 5, 2026-09-23 (note 2077): general framing; no survey of other countries'
% inventory systems (too prescriptive). CBM-CFS3/Century/RothC primary reading dropped.
\gap{To be written, in general terms and without naming other national inventories.
Content: what follows for any application that starts a soil carbon model at steady
state; reading a calibrated input multiplier as a physical quantity that can be
compared with independent budgets; and the difference between skill at the scale of a
national mean and skill between individual plots.}

\subsection{A national map of soil carbon change as a by-product}
% ROUND 5, 2026-09-23 (note 2162): moved here from the Conclusions.
% ROUND 5, 2026-09-23: from Reporting/NextgenC_report/ZENODO_description.md (maps rebuilt
% on run 20260831, 2026-09-03). Nugget 85-90% of the variogram; 2015-2024 ensemble delta
% ~87% of cells gaining, mean +0.18 tC/ha/yr.
The calibrated models also yield a gridded product. For each model and for the ensemble
of the six, weighted equally, the posterior mean and standard deviation of the soil
carbon stock of each plot and year were interpolated by ordinary kriging on the
logarithmic scale onto a $2$ km grid of the mineral-soil forest land of Finland for
1985--2024, with cells dominated by peatland or water masked. The standard deviation of
each map combines the uncertainty of the interpolation with the posterior uncertainty of
the model and, for the ensemble, the spread between the six models. Over 2015--2024 the
ensemble gains carbon in about $87$ per cent of the cells, at a mean rate of $0.18$ tonnes
of carbon per hectare per year. The maps are a deliverable and not a finding: the plots
are on average about $27$ km apart, and at that distance the stock of neighbouring plots
is almost unrelated, so the maps reproduce the national field with local deviations
anchored on the plots, and the detail between plots is interpolation rather than
measurement. The maps are published as an open dataset \gap{Zenodo DOI of the map
deposit, once deposited}.
\gap{Before publication, check the region of carbon loss in the far south of the delta
map (about $-0.5$ tonnes per hectare per year): a real feature of the plots, or an
artefact of kriging near the edge of the grid?}

% =============================================================================
\section{Limitations}

% ROUND 5, 2026-09-23: first draft. No reference to our own development history.
The observed trend rests on three campaigns that were not carried out in the same way.
The first campaign spread its fieldwork over a decade, recorded the organic layer
without its litter and moss, and predicted the carbon content of most of its mineral
layers from the loss of mass on ignition. We standardised what the measurements allow,
but a shared offset between campaign levels cannot be separated from a change of the
stock: the rate over the full observed window survives an offset of about $9$ per cent,
whereas the rate over the last interval, on which the apparent slowing of the
accumulation rests, is removed by an offset of about $2$ per cent.

The history of the soil before the first observation is represented coarsely. The path
of the litter input from 1917 to 1985 is one national shape, taken from a reconstruction
whose tree-derived terms cover all forest land while the area used to divide them covers
upland soils only, so that the rise of the input before 1985 is, if anything, overstated.
The starting state is set by one national parameter, whereas the history of exploitation
that left the soil out of balance differed between regions and between plots; the models
therefore cannot represent why individual plots differ, which is part of the reason for
their lack of skill between plots. The starting state also has a single degree of freedom: the soil of 1917
is assumed to be in balance with the litter input of 1917, so that a soil depleted below
that balance can be expressed only as a lower input, and the projected sink depends
strongly on how that deficit is estimated (Section~\ref{sec:headroom}). The stock of 1917 implied by each calibration is a
consequence of the calibration and is not constrained by any observation
\gap{report the implied 1917 stocks on the current run, from
doublechecks/init\_state\_plausibility.R, and compare them with what is known of boreal
forest soils; the floor for that comparison still has to be sourced}.

The error model rests on design values. The total spread of the error, $0.800$ on the
logarithmic scale, is declared rather than estimated, and the spreads of the campaign
offsets are prescribed rather than measured.
\gap{Decide whether to test the sensitivity of the main results to the total spread of
the error (one additional set of calibrations), or to state here that it was not tested.}

The effect of estimating the starting state has not been isolated. The comparison with the
published parameterisations mixes the starting state with every other difference between
this calibration and the published ones, from the data to the error model, and a
calibration started from a soil in balance with the present input, which would isolate
the effect, was not carried out.
\gap{This paragraph limits what the title and the Conclusions may claim. Settle together
with the title: run the balanced-start calibration, or phrase every causal statement
about the starting state as a hypothesis.}

The projections are orders of magnitude and not forecasts. The litter input is held at
its value in 2024 and the climate of the last $20$ observed years is repeated, so the
projections contain neither a change of climate nor a change of forest growth or
management, and the projected sink is the relaxation of the soil towards balance under
present conditions.

% =============================================================================
\section{Conclusions}

% ROUND 5, 2026-09-23 (notes 2107, 2114, 2130, 2139, 2149, 2158): rewritten. No
% "reach the stocks"; no claim about what a balanced start cannot do (the
% counterfactual has not been run); "suggest", not "require"; the projection effect
% named per model; saturation stated plainly from Section sec:headroom; the study
% presented as a baseline for model development.
A transient starting state, estimated rather than assumed, allows all six models to
represent the soil carbon measured in the three campaigns and to bracket the observed
rate of change over the full observed window. The fit to these data cannot separate
the six structures, and the models differ where the data do not constrain them: in
the shape of their trajectories, and therefore in their extrapolation.

Under the assumptions of this study, reproducing the observed trajectory places the
models at a larger litter input than the product supplies and, in two of the three
versions of Yasso, at a faster turnover than the published parameterisations. The two
are one displacement along a trade-off that the soil measurements cannot resolve, and
moving back towards the published transit times costs little in fit, so the result
suggests a direction rather than a value. The preference for a larger litter input
follows the measurements more than the prior. The most plausible reading of the
displacement is not that the models are wrong but that the litter input is incomplete,
and the components missing from it by construction --- carbon released by living
roots, understorey vegetation, mycorrhizal residues and early-succession vegetation ---
are of the right kind, and together of a size that brackets the gap.

Where the calibration lands on the trade-off changes the projection of Yasso07 and not
those of Yasso15 and Yasso20. A model that applies one climate modifier to every pool
ties its turnover to its sensitivity to temperature, so that moving along the
trade-off moves the projection; models that apply separate modifiers to separate pools
are insensitive to the same movement, and are therefore more robust to a calibration
choice that the present data cannot settle.

Under a constant litter input, the soil carbon sink of Finnish upland forests is
saturating. The models place the soil $26$ to $42$ per cent below its state of balance,
in agreement with the $38$ per cent estimated by \citet{Liski2006}, yet the sink they
project falls by a third to a half within $60$ years, because most of the remaining distance
lies in the slowest pool and is realised only over centuries. How much sink remains
depends on how far below balance the soil was at the start, which the calibration infers
only indirectly: the soil of 1917 appears more depleted than its litter input alone would
imply, and a soil assumed to start in balance with its present input would have almost
nothing left to gain. A larger litter input
buys additional storage, $6$ to $10$ tonnes of carbon per hectare by 2050 for an
increase of $20$ per cent, but the gain is finite, largest in the first decades, and
lasts only as long as the larger input is maintained. Both statements come from
projections without a change of climate and describe orders of magnitude, not
forecasts.

About half of the variation between individual plots that none of the models
represents can be recovered from site variables that already exist. This study is
therefore a baseline for further model development: it shows where the current models
fall short, in the differences between individual plots, and which information, on the
productivity of the stand and on the physical properties of the soil, a next model
would need to include.

% =============================================================================
\section*{Acknowledgements}
% Items collected as they arise; the section is written last. O.-P. Tikkasalo (the
% weather extraction) is expected to be a coauthor and is therefore NOT listed here.
\gap{To be written. Collected so far: LUKE Data Support (weather-data provenance,
2026-09); the gridded daily weather observations are open data of the Finnish
Meteorological Institute; the soil carbon measurements and their provenance
(H.~Ilvesniemi's workbook) go here and in the data statement, not in the references.}

\section*{Data and code availability}
\gap{To be written last. The litter product is on Zenodo; the calibration and
analysis code are in the project repository and will be archived on Zenodo with a DOI
at submission, so that the version used is fixed.}

\bibliographystyle{plainnat}
\bibliography{library}
\end{document}
